\documentclass[11pt]{amsart}
\usepackage{lmodern}
\usepackage[expansion=false]{microtype}
\usepackage[a4paper,textwidth=152mm,textheight=242mm,centering]{geometry}
\usepackage{amsmath,amssymb,mathtools,booktabs,enumitem}
\usepackage[colorlinks,linkcolor=blue!50!black,citecolor=blue!50!black,urlcolor=blue!50!black]{hyperref}
\usepackage{xcolor}
\setlength{\emergencystretch}{1.4em}
\numberwithin{equation}{section}
\newtheorem{theorem}{Theorem}[section]
\newtheorem{lemma}[theorem]{Lemma}
\newtheorem{proposition}[theorem]{Proposition}
\newtheorem{corollary}[theorem]{Corollary}
\theoremstyle{definition}
\newtheorem{definition}[theorem]{Definition}
\newtheorem{example}[theorem]{Example}
\newtheorem{remark}[theorem]{Remark}
\theoremstyle{plain}
\newtheorem{mainthm}{Theorem}
\newtheorem{maincor}[mainthm]{Corollary}
\renewcommand{\themainthm}{\Alph{mainthm}}
\DeclareMathOperator{\lcm}{lcm}
\DeclareMathOperator{\Tr}{Tr}
\DeclareMathOperator{\FPdim}{FPdim}
\DeclareMathOperator{\Irr}{Irr}
\DeclareMathOperator{\Gal}{Gal}
\newcommand{\QQ}{\mathbb Q}
\newcommand{\ZZ}{\mathbb Z}
\newcommand{\C}{\mathcal C}
\allowdisplaybreaks[2]

\title[Rank bounds for simple integral modular data]{Rank bounds for simple integral modular data}
\author{Sebastian Burciu}
\address[S. Burciu]{Simion Stoilow Institute of Mathematics of the Romanian Academy, Bucharest, Romania}
\email{sebastian.burciu@imar.ro}
\author{Jingcheng Dong}
\address[J. Dong]{College of Mathematics and Statistics, Nanjing University of Information Science and Technology, Nanjing, China}
\email{jcdong@nuist.edu.cn}
\author{S\'ebastien Palcoux}
\address[S. Palcoux]{Beijing Institute of Mathematical Sciences and Applications, Beijing, China}
\email{sebastien.palcoux@gmail.com}
\date{}
\hypersetup{pdftitle={Rank bounds for simple integral modular data},pdfauthor={Sebastian Burciu, Jingcheng Dong and S\'ebastien Palcoux}}
\begin{document}
\begin{abstract}
We prove that a non-pointed based-simple integral fusion ring with symmetric Drinfeld duality, cyclotomic modular data and rotation certificates has rank at least $45$, and at least $217$ in odd dimension. The proof bounds the possible degrees and their multiplicities using character orthogonality and the primary components of the twists.
\end{abstract}
\maketitle

\section{The rank problem}
The five irreducible degrees of $A_5$ are $1,3,3,4,5$. Its representation ring is based-simple \cite[\S5]{LPR}, but $3^2\nmid60$ prevents its character table from having the symmetry of integral modular data. How large must the rank be when a fusion ring has this symmetry and satisfies the modular identities for twists and rotations?

\smallskip\noindent
A commutative fusion ring $R$ is a commutative unital ring with finite
basis $(x_i)_{i\in I}$ over $\ZZ$, unit $x_0$ (the vacuum), nonnegative integral
structure constants $N_{ij}^k$, and a basis-preserving ring involution
$x_i\mapsto x_{i^*}$ such that $N_{ij}^0=\delta_{i,j^*}$.
A character is a unital homomorphism $R\to\mathbb C$. Write $d_i$ for the value at $x_i$ of the unique positive character,
$r=|I|$ for its rank, and $D=\sum_i d_i^2$ for its global dimension.
The ring is \emph{integral} if all $d_i$ are integers, \emph{pointed}
if all equal one, and \emph{based-simple} if every nonunit $x_i$
generates the full basis under duality and passage to constituents of
products. A \emph{symmetric algebraic Drinfeld map} is a bijection
from the probability basis $c_i=x_i/d_i$ to the complex characters
$\phi_i$ satisfying $\phi_i(c_j)=\phi_j(c_i)$. Put
$\chi_i(j)=\phi_i(x_j)$, $S_{ij}=d_i\chi_i(j)$, and $s=S/\sqrt D$.

\begin{mainthm}\label{alg:rank}
Let $R$ be a non-pointed based-simple commutative integral fusion ring whose probability basis admits a symmetric algebraic Drinfeld map. Suppose that $R$ carries rotation-admissible cyclotomic modular data as in Definition~\ref{alg:mod-data}. Then $\operatorname{rk}(R)\ge45$. If $D$ is odd, $\operatorname{rk}(R)\ge217$.
\end{mainthm}

\begin{definition}\label{alg:mod-data}
\emph{Cyclotomic modular data} consist of roots of unity $\theta_i$
with $\theta_0=1$ and a root of unity $\gamma$ of two-power order.
Set $T=\operatorname{diag}(\theta_i)$,
$N=\operatorname{ord}(T)$ and $L=\operatorname{ord}(T/\gamma)$, and write $\zeta_n=e^{2\pi i/n}$.
The primes dividing $N$ must be precisely those dividing $D$, and
the assignments
\[
 W=\left(\begin{smallmatrix}0&-1\\1&0\end{smallmatrix}\right)
 \longmapsto s,\qquad
 U=\left(\begin{smallmatrix}1&1\\0&1\end{smallmatrix}\right)
 \longmapsto T/\gamma
\]
must define a representation $\rho$ of $\mathrm{SL}_2(\ZZ/L)$ over
$\QQ(\zeta_L)$ with
\begin{equation}\label{alg:covariance}
 \sigma_a\rho(g)=\rho(D_a gD_a^{-1}),\qquad
 D_a=\operatorname{diag}(a,1),\quad \sigma_a(\zeta_L)=\zeta_L^a
 \quad (a\in(\ZZ/L)^\times).
\end{equation}
For $m\ge1$ define the linear functional
\begin{equation}\label{alg:rotation}
 L_m^i(x_c)=\sum_{a,b}s_{ia}s_{ib}(\theta_a/\theta_b)^m N_{ab}^c .
\end{equation}
A \emph{rotation certificate} for $(i,m,c)$ is a finite multiset
of $m$th roots of unity whose $t$th powers sum to $L_{m/t}^i(x_c^t)$
for each positive divisor $t$ of $m$. We call the data
\emph{rotation-admissible} if every $(i,m,c)$ has a certificate.
\end{definition}

Equip an integral modular fusion category with its canonical positive spherical structure \cite[Propositions~8.23--8.24]{ENOannals}. Its normalized $S$-matrix gives the character pairing, and generalized Frobenius--Schur rotations give the required power traces \cite{DLNCongruence,NSind}. The existence of a non-pointed simple integral modular fusion category is open \cite[Question~1.3 and Theorem~5.8]{LPR}. For categories the previously known lower bounds are $32$ in general and $57$ in odd dimension \cite[Theorem~22.2]{BDPcore}.

Two features make a rank bound possible. Write $D=qM^2$ with $q$ square-free. The formal codegree $D/d_i^2$ is an integer, so $b_i=M/d_i$ is an integer and $\sum_i b_i^{-2}=q$ \cite[Lemma~2.3]{HR}, \cite[\S4]{ABPP}. At small rank, some of these denominators must be small. At a prime $p\mid D$, reduction of the characters partitions the labels into blocks, while the $p$-part of each twist prescribes orthogonality relations within those blocks. They constrain the multiplicities of the largest degrees and the degrees on which their character rows can be nonzero. In odd dimension duality pairs all nonunits, and the same relations constrain whole pairs at once. The finite counts in Sections~\ref{sec:odd-new} and~\ref{sec:verify} complete the proof.

\section{Character arithmetic}
Throughout the proof, $R$ and its modular data satisfy the hypotheses of Theorem~\ref{alg:rank}.

\begin{lemma}[The normalized character table]
The matrix $s$ is symmetric and unitary, $s^2$ is the duality
permutation $C$, and
\begin{equation}
 \sum_j\chi_i(j)\overline{\chi_k(j)}=\delta_{ik}f_i,\qquad
 f_i=D/d_i^2\in\ZZ,\qquad S_{ij}/\lcm(d_i,d_j)\in\overline{\mathbb Z} .
\end{equation}
The character field $K_i=\QQ(\chi_i(j):j\in I)$ has conductor dividing
$f_i$. Galois permutes the labels, preserving degree, with
$\sigma S_{ij}=S_{\sigma i,j}=S_{i,\sigma j}$.
\end{lemma}
\begin{proof}
The fusion matrices are commuting normal matrices, their adjoints
being given by duality. Their distinct character rows are orthogonal.
Symmetry gives $SS^*=S^*S=\operatorname{diag}(d_i^2f_i)$, where
$f_i=\sum_j|\chi_i(j)|^2$. This diagonal commutes with $S$ and its
vacuum column has no zero, so it is $DI$. Also $\bar S=CS$.
The entries $\chi_i(j)$ are algebraic integers. This gives integrality
of the rational $f_i$, and B\'ezout applied to $S_{ij}/d_i,S_{ij}/d_j$
gives the lcm assertion. Galois permutes the characters. Their
codegrees, hence their degrees, are fixed. Symmetry then gives
covariance in both indices. The row-conductor theorem
\cite[Theorem~2.9]{BDPcore} applies to these symmetric integral data.
\end{proof}

As $\rho(-I)=C$, dual labels have the same twist. The identity
$WUW=U^{-1}WU^{-1}$ and its vacuum entry give
\[
 sTs=\xi T^{-1}sT^{-1},\quad \xi=\gamma^3=p^+/\sqrt D,\quad p^+=\sum_i d_i^2\theta_i .
\]
For $v=(d_i\theta_i)_i$, its vacuum column gives
$(s^{-1}v)_l=\sqrt D\,\xi s_{0l}/\theta_l$. Diagonalizing fusion
and using the inverse of the same identity gives the balancing identity
\[
 \sum_k N_{i^*j}^k d_k\theta_k
 =\sqrt D\,\xi(sT^{-1}s^{-1})_{ji}=\theta_i\theta_j S_{ij}.
\]
Thus $S_{ij}\in\QQ(\zeta_N)$. Write $h_a=\rho(\operatorname{diag}(a,a^{-1}))$.
By \eqref{alg:covariance}, $h_as=\sigma_a(s)$. Row covariance makes
$h_a$ a signed permutation fixing the vacuum line. Conjugacy of $U$ to $U^{a^2}$ by
the diagonal matrix therefore gives $\gamma^{a^2-1}=1$ for every
unit $a$ modulo $L$. At a prime of $L$ outside $D$, both generating
unipotents of the local factor act by scalars.

\begin{lemma}[Arithmetic of rotation traces]\label{alg:rotation-arithmetic}
Every nonnegative integral combination $V$ of basis elements has
rotation certificates. All their power traces belong to $\QQ(\zeta_N)$.
For $p\mid m$ and $y$ with algebraic-integer coefficients,
$L_m^i(y)^p\equiv L_{m/p}^i(y^p)$ at every place over $p$.
For $m=N_{p'}:=N/p^{v_p(N)}$, $L_m^0(x_c)\in\ZZ$.
\end{lemma}
\begin{proof}
Take the summands of $V$ as letters, with distinct copies distinguished.
For each cyclic orbit of length-$m$ words of least period $b$, let $Y$
be the product along a period. Induce its certificate from $C_{m/b}$
to $C_m$. This exists by hypothesis for $b=1$ and by induction on $m$
otherwise. Its trace at $t\mid m$ is zero unless $b\mid t$, when it is
$bL_{m/t}^i(Y^{t/b})$. Summing counts each length-$t$ word once.
Divisor-power traces are in $\QQ(\zeta_N)$ by \eqref{alg:rotation}
and $s_{ia}s_{ib}=S_{ia}S_{ib}/D$. Other power traces are their
Galois conjugates. Reduction of the roots gives the congruence on the
basis, hence by Frobenius on all $y$. For the last assertion the
integral trace lies in both $\QQ(\zeta_m)$ and $\QQ(\zeta_{p^{v_p(N)}})$,
the latter by the rational vacuum weights in \eqref{alg:rotation}.
\end{proof}

\begin{lemma}[Vanishing and divisibility]\label{alg:odd-parity}
Only the unit has degree one. The dimension is square-full, with
at least three prime divisors, and every nonunit degree has at least
two prime divisors. For nonunits $i,j$ one has $|S_{ij}|<d_id_j$,
with $S_{ij}=0$ when $(d_i,d_j)=1$. After deleting any single nonunit
Galois orbit the remaining nonunit degrees have gcd one.
If $D$ is odd, every nonunit is paired with a distinct dual, its row
field is nonreal, and $D\equiv r\pmod{16}$.
\end{lemma}
\begin{proof}
The invertibles span a proper based subring, so reduce to the unit.
This property is called \emph{perfection}. By Burnside vanishing \cite[Theorem~1.2]{BPBurnside}, the product
$\prod_i x_i x_{i^*}$ is supported only on the degree character.
Its integral unit coefficient is $\prod_i d_i^2/D$. Since $d_i^2\mid D$,
every prime divisor of $D$ occurs at least twice. For fixed nonunit $i$ the equality
labels $|S_{ij}|=d_id_j$ span a based subring
\cite[Lemma~2.35]{BPBurnside}. Since $f_i<D$ it is proper.
For coprime degrees every conjugate of the algebraic integer
$S_{ij}/(d_id_j)$ therefore has modulus less than one. Vacuum
orthogonality gives $\sum_{j\ne0}d_j\chi_i(j)=-1$. If $d_i$ were a
prime power, every nonzero term would be divisible by that prime.
The same identity gives gcd one for all nonunit degrees, hence at
least three dimension primes. By symmetry every nonunit column has a zero. A row vanishing
there vanishes on its entire Galois orbit, giving the deletion claim.

For odd $D$, the unit coefficient
$\tau(y)=D^{-1}\sum_i d_i^2\phi_i(y)$ reduces above two to the sum
of the characters. Thus $\tau(x_i^2)\equiv\tau(x_i)^2\pmod2$,
or $\delta_{i,i^*}=\delta_{i,0}$. Conjugation acts by duality, so
nonunit row fields are nonreal. In each dual pair $2(d_i^2-1)$ is
divisible by sixteen.
\end{proof}

\paragraph{Ordered denominators.}
Write $D=qM^2$ with $q$ square-free. Then $d_i\mid M$. Put
$b_i=M/d_i$ and order the nonunit denominators increasingly, with
repetition. For a prime $p$ write $M_p=p^{v_p(M)}$. A $p$-free
degree is prime to $p$, and a pure $p$-denominator is a power of $p$.
Primes dividing $q$ are called ramified, and the other primes of $D$
unramified. Dividing the mass $\sum_{i\in J}d_i^2$ of a set of
labels by $M^2$ gives its reciprocal-square mass $\sum_{i\in J}b_i^{-2}$.
In odd dimension $r\equiv q\pmod8$. We have
\begin{equation}
 q=M^{-2}+\sum_{i=1}^{r-1}b_i^{-2},\qquad qb_1^2<r .
\end{equation}
After a prefix of length $j\ge1$ write $L=\lcm(b_1,\ldots,b_j)$,
$E=L^2(q-\sum_{i\le j}b_i^{-2})$, $h=r-1-j$. For the next $b$,
\begin{equation}\label{eq:interval}
 b\ge b_j,\qquad L^2<Eb^2\le(h+1/36)L^2 .
\end{equation}
Indeed $M/b\ge6$. For $g=(L,b)$ the new pair is
$((b/g)L,(b/g)^2 E-(L/g)^2)$. With two degrees $d\le e$ remaining,
write $M=kL$. Each prefix degree is a multiple of $k$. Deleting
the orbit of either remaining label shows $(k,de)=1$. Thus
$d,e\mid L$, $Ek^2=1+d^2+e^2$, and $kL/e\ge b_j$.
Reduction modulo four and five gives $(k,10)=1$ and then
$E\not\equiv0,4,7\pmod8$. With one degree left, deletion gives
$k=1$, $E=1+d^2$.

\paragraph{Galois orbits and projection.}
\begin{lemma}[Projection onto an orbit]\label{result:field-rank}
Let $O$ be a Galois orbit of columns and $Y\in O$. Divide orthogonal
rows, equivariant in the column and having algebraic-integer entries
nonzero at $Y$, into families whose values at $Y$ generate fields $F_j$.
Then $[\prod F_j:\QQ]\le |O|$ (the product denotes the compositum) and
\[
 |O|\sum_v\|v\|^{-2}\le1+\sum_j([F_j:\QQ]-1).
\]
If these rows are orthogonal to the degree row $d$, one may add
$|O|/f_Y$ on the left.
\end{lemma}
\begin{proof}
The stabilizer of $Y$ fixes the compositum. On $O$ each family lies
in the function space on the embeddings of $F_j$. These spaces share
the constants. The restricted matrix of normalized rows is a contraction,
so its squared Hilbert--Schmidt norm is at most this dimension.
Each entry at $Y$ has mean squared conjugate modulus at least one
by the norm inequality, giving the lower bound. Adjoining $d/\sqrt D$
adds a constant row on $O$. We shall also use the projection inequality
\begin{equation}
 I-\sum_v vv^*/\|v\|^2\succeq dd^*/D
\end{equation}
for orthogonal rows perpendicular to $d$, regarded here as columns.
\end{proof}

\paragraph{Cyclotomic estimates.}
Write $\langle z\rangle$ for the mean of the Galois conjugates of $z$
and $C(z)=\langle |z|^2\rangle$. We use $v_p(p)=1$.
A nonzero cyclotomic integer satisfies $C(z)\ge1$. Consequently
$C(z)\ge|\Re\langle z\rangle|$ by Cauchy--Schwarz.

\begin{lemma}\label{result:coupled-gap}
Let $A$ be a sum of an even number of roots of two-power order. If $A\ne0$, then $C(A)\ge2$.
\begin{enumerate}[label=\textup{(\roman*)},nosep]
\item If $B$ is a sum of at most $k$ roots, $u>k/2$ is rational,
and $(A,B)\ne(0,0)$, then $C(uA+B)\ge\min\{1,2(u-k/2)^2\}$.
\item Suppose the roots in $B$ have orders supported on $\{2,3\}$.
If $B$ has at most eight terms and $9A+4B$ is a rational linear
combination of roots with total absolute coefficient less than two,
then $A=B=0$.
\item If $B$ as in \textup{(ii)} has at most five terms, $E$ has zero
or two dyadic terms, and the same coefficient bound for $36A+16B+9E$
is less than six, then $(A,B,E)$ is zero or $(0,-\eta,2\eta)$ for
a dyadic root $\eta$.
\end{enumerate}
\end{lemma}
\begin{proof}
Normalized trace onto a dyadic field containing all dyadic parts
multiplies each dyadic part by the rational mean of its odd part.
In its orthonormal power basis the projected $B$ has coefficient
$\ell^1$-norm at most $k$, while the integral vector $a$ of $A$ has
even coordinate sum, so $C(A)=\|a\|^2\ge2$ for $a\ne0$. On $m\ge2$ nonzero coordinates of $a$ the
squared norm of the projected $uA+B$ is at least
$(um-k)^2/m\ge2(u-k/2)^2$. For $m=1$ its bound is $(2u-k)^2$.
For $m=0$ use $C(B)\ge1$.

In (ii)--(iii) write $Z$ for the displayed combination. If $B$ is
outside the dyadic field, its squared distance from it is at least
$3/4$. Indeed an orthogonal component above the cubic-root extension
has norm at least one, while $P+Q\zeta_3$ with dyadic $P,Q$, $Q\ne0$,
has squared distance $3C(Q)/4$ \cite[Lemma~5.2.1]{CMSCyclotomic}.
This gives $C(Z)\ge12$ or $192$, respectively, exceeding the squared
coefficient bound. Thus $B,Z$ have integer coordinate vectors $b,z$
in the dyadic basis. In (ii) $\|z\|_1<2$ and even coordinate sum
give $z=0$. Now $B=-9\delta$, $A=4\delta$ and all conjugates
of the integer $\delta$ have modulus at most $8/9$.

In (iii) $\|b\|_1\le5$ and $\|z\|_1<6$. At $e_j=z_j=0$ the
relation forces $a_j=b_j=0$. For $e=0$, a nonzero $z$ has one
entry $\pm4$, forcing the sole nonzero $a_j$ odd. If $e$ has
one entry, change its sign to $2$. There $z_j=\pm2$ and all other
entries of $z$ vanish. Thus $a_j$ is even and $9a_j+4b_j=-4$
or $-5$, forcing $a_j=0$, $b_j=-1$. For two entries $\pm1$ of
$e$, the absolute $z$-values there are $(1,1)$ or $(1,3)$ up to order.
Modulo $36$, $16b_j=z_j-9e_j$ forces $|b_j|\ge4$ for $|z_j|=1$
and $|b_j|\ge3$ for $|z_j|=3$, exceeding five in total.
\end{proof}

\paragraph{Principal blocks.}
At a place above $p\mid D$, labels form a $p$-block when their
characters have the same reduction on $R$. The partition is independent
of the place by covariance on columns. Over the completed valuation
ring $\mathcal O$ of a splitting field the block idempotent
$e_B=\sum_{j\in B}e_j$ is integral by idempotent lifting
\cite[Tag 04GH]{Stacks}. The character transform identifies $e_j$ with the $j$th coordinate vector, also denoted $v_j$.
Consequently
\begin{equation}\label{eq:block-trace}
 \sum_{j\in B} d_j^2\chi_j(y)=D\tau(e_By)\in D\mathcal O
 \qquad(y\in R\otimes\mathcal O).
\end{equation}
We use $\chi_j(y)$ also for the linear extension of $\chi_j$.

\begin{lemma}\label{block:height}
The principal block $P=P_p$ contains the vacuum and is Galois-stable.
One has
\begin{equation}\label{eq:principal-trace}
 p^{v_p(f_i)}\mid\sum_{j\in P}|\chi_i(j)|^2,\qquad
 p^{v_p(D/d_i)}\mid\sum_{j\in P}d_j\chi_i(j)
\end{equation}
in $\ZZ$. A row of $p$-power codegree is supported in $P$, with
support equal to $P$ if $p\nmid d_i$. If $p^2\mid D$, then
$p^2\mid f_i$ for every $i\in P$.
\end{lemma}
\begin{proof}
The norm divisibility and height bound are the principal-block
theorems \cite[Theorem~2.11 and Corollary~4.1]{BDPcore} for the
symmetric order $R$, whose unit-coefficient pairing makes the fusion
basis orthonormal against its dual. The norm contains the vacuum
contribution one, so for $p$-power $f_i$ it exhausts the row norm.
If $p\nmid d_i$, then for $j\in P$ the entry $\chi_j(i)\equiv d_i$
is nonzero. The second sum equals $(D/d_i)\tau(e_Px_i)$ by symmetry
and is a rational integer by Galois invariance.
\end{proof}

\section{Blocks and primary phases}
The Chinese remainder theorem factors $\mathrm{SL}_2(\ZZ/L)$ as the product of its prime-power groups. For $p\mid D$ this gives local Weyl, unipotent
and torus operators $W_p,U_p,H_p(a)$, the last corresponding to
$\operatorname{diag}(a,a^{-1})$. The $p$-phase of $i$ is the
$p$-part $\theta_{i,p}$ of its twist. We say that row $i$ meets label $j$ if $S_{ij}\ne0$. A label is active at $p$ if
this phase differs from one. A phase fiber fixes any specified
primary phases. Its mass is the sum of the $d_i^2$ in it. Its
projector commutes with every unspecified local factor.
For a multiplier $a$ supported at $p$, covariance gives
$H_p(a)s=\sigma_a(s)$. The torus acts by the Galois permutation
(up to the common sign $\sqrt D/\sigma_a(\sqrt D)$), and raises
normalized phases to square-unit powers.

\begin{lemma}[Local block rules]\label{block:admissible}
Every full $p$-block spans a local invariant subspace. A row meets
a $p$-active label precisely when $p\mid f_i$. Its coordinate line
is $W_p$-invariant precisely when $p\nmid f_i$. Activity of $i$
requires $p\mid f_i$, and for odd $p$ its $p$-phase order divides
$f_i$. A label of $p$-power codegree belongs to the principal block
$P_p$. Then every $p$-free degree belongs to $P_p$ and has
$p$-power twist. A local invariant coordinate plane on which the
Weyl moves one coordinate line occurs only at two, with equal degrees
and phase ratio of order four.
\end{lemma}
\begin{proof}
For $A_j=sT^js^{-1}$, substituting the primitive idempotents in
\eqref{alg:rotation} gives
\begin{equation}\label{eq:rotation-entry}
 L_j^i(e_z)=|(A_j)_{iz}|^2,\qquad
 (L_j^i(y))^p\equiv L_{j/p}^i(y^p)\pmod{\mathfrak p}\quad(p\mid j),
\end{equation}
the latter for integral coefficients by the rotation certificates.
Here and below we use the ordinary-table arguments of
\cite[\S\S3,6,19]{BDPcore}. To check their hypotheses,
Lemma~\ref{alg:rotation-arithmetic} realizes every divisor trace on
a nonnegative integral element as a power trace of a single multiset
of roots, with all power traces in $\QQ(\zeta_N)$. This is precisely
the rotation input in \cite[Lemma~3.1]{BDPcore}: Frobenius gives the
congruence above, and Galois invariance of multiplicities gives its
higher-power congruence (3.1). Its Theorem~3.2 uses only the diagonal
insertions $i,i^*$, so applies to our $L^i_j$ and gives block
orthogonality. The local action, line and order criteria and the plane
rule follow by the proofs of its Lemmas~3.3--3.4, Theorem~3.6 and
Lemma~19.5. Their remaining inputs are congruence covariance, conductor
divisibility, $\operatorname{rad}N=\operatorname{rad}D$ and the
positive rational vacuum ratios. For prime-power codegree the
complementary Weyls preserve the line. Its $s$-column therefore
lies in its $p$-block. The vacuum entry makes the block principal.
All $p$-free degrees meet this row by reduction, so belong to its
block and have $p$-primary twist by the row criterion.
\end{proof}

\begin{lemma}[Principal support]\label{block:witness}
The principal block $P=P_p$ contains at least three distinct nonunit
degrees. They share a prime in pairs, have gcd one, and every one
meets each $p$-free nonunit degree in a common prime. For any
partition of the primes of $M$ into two or more nonempty parts,
at least three distinct nonunit denominators meet two parts, and
their lcm is $M$.

The $p$-free nonunits of $P$ occupy at least two Galois orbits.
At $p=3$ in odd dimension there are at least eight such labels.
\end{lemma}
\begin{proof}
Orthogonality of an active row to the vacuum on $P$ gives
$\sum_{0\ne Y\in P}d_Y\chi_X(Y)=-1$. The degree gcd follows.
A principal character is nonzero on every $p$-free degree by
reduction. Coprime vanishing then gives the asserted pairwise
intersections and hence at least three distinct degrees. On $P$,
$p\mid b_Y$ by Lemma~\ref{block:height}. Given a prime $t\mid M$,
choose $p$ in another part. Some nonunit in $P_p$ has $t$-free
degree, proving the lcm assertion. If only two crossing degrees
occur they are coprime. For $p$ in one part $J$ all other nonunit
degrees in $P_p$ are multiples of the part of $M$ outside $J$.
Thus exactly one crossing degree belongs to $P_p$, and it is
supported on $J$. Repeating in another part gives a coprime
$p$-free degree, contradicting principality.

For every row $X$, the principal trace
$\sum_{Y\in P}d_Y\chi_X(Y)=D\tau(e_Px_X)/d_X$ is zero modulo $p$.
Its vacuum term is one, so $X$ meets a $p$-free nonunit of $P$.
Burnside vanishing supplies a row vanishing on any one chosen
nonunit orbit, proving the two-orbit assertion. In odd dimension
the count at three is even and at least four, and is $2\bmod3$
by \eqref{eq:block-trace} at the unit.
\end{proof}

\begin{lemma}[Counting phases]\label{block:moments}
Fix an odd prime $p$. In a complementary phase fiber, also after
restriction to a full $p$-block, primitive $p$th roots have equal
mass when $v_p(D)$ is even. For odd valuation the mass at order $p$ is
$(p-1)$ times the neutral mass. At fixed degree and complementary
phases an order-$p^e$ orbit has $p^{e-1}h_p$ distinct phases, where
$h_p=(p-1)/2$.
Put $A=\{t\mid q:t\text{ odd prime}\}$,
$C=\prod_{t\in A}(h_t+1)$, and
\[
 B(q)= C\max\bigl(\{2^{\mathbf1_{2\mid q}}\}
               \cup\{(h_t+2)/(h_t+1):t\in A\}\bigr).
\]
Every nonempty fiber fixing odd-primary phases outside $A$ has
at least $B(q)$ labels. For a set $J$ of odd primes of positive even
valuation, at least $\sum_{p\in J}(p-1)$ nontrivial $J$-phase tuples
occur. At every prime $p\mid D$, including two, the $p$-active mass $V_p$ satisfies
$2pxV_p\ge(p-1)D$ for every active degree $x$ and for the degree
of every based generator.
\end{lemma}
\begin{proof}
Apply the ordinary-table slice formulas of
\cite[Propositions~3.9, 19.1, 19.8 and \S19.2]{BDPcore}. Their
inputs are block invariance and the Bruhat identity on a commuting
fiber projector. The torus supplies the indicated square orbits. For
the last estimate the integer $L^0_{N_{p'}}(X)$ is strictly less
than $x$ if a fusion edge changes $p$-phase. This holds at the unit
for an active label and by connectedness for a generator. Averaging
the roots with weights $N_{ab}^X d_a d_b/D$ bounds the deficit by
$2pxV_p/((p-1)D)$, as in \cite[Proposition~3.8]{BDPcore}.
\end{proof}
A packet consists of $h_p$ equal-degree labels on a square orbit.
Degree $d$ of multiplicity $m_d$ supplies at most $n_d=\lfloor m_d/h_p\rfloor$
packets. At even valuation the differences $k_d$ of the two primitive
packet counts satisfy $|k_d|\le n_d$ and $\sum k_dd^2=0$. If
$n_d-|k_d|$ is odd and either $<p$ or $p^2\nmid f_d$, at least
one place is unused: further primitive packets come in pairs and
a higher orbit costs $p$ places. Two packets alone therefore have
equal degree. At $2\mid q$ the plane rule gives
$\sum_d(n_d+\lfloor n_d/2\rfloor)\ge6$
\cite[Lemma~19.6]{BDPcore}.

\begin{lemma}[Local tensor factors]\label{result:spherical}
Call a prime of $M$ outside $q$ allowed. For a set $P$ of allowed
primes, let $F$ run over the maximal subsets occurring together in
a twist. For each $K\subseteq F$ one may simultaneously reserve
$C\prod_{p\in K}(p-1)$ labels with denominator divisible by
$\prod_{p\in F}p$ and active within $P$ exactly on $K$. In odd
dimension these are reservations of duality orbits. If a label is
active on $J\subset P$ and $\ell\in P\setminus J$ divides its
denominator, a twist is active throughout $J\cup\{\ell\}$.
\end{lemma}
\begin{proof}
Decompose the vacuum cyclic module into irreducible tensor products.
Each factor contains a vector $v$ with the vacuum Borel character.
At odd $t\mid q$ the vacuum Borel character is quadratic.
Thus the local factor has a neutral vector and a nontrivial
unipotent eigenvalue (the conjugates of $U_t$ generate). Its
square-unit translates give at least $(t+1)/2$ dimensions. At an allowed odd $p$ this character is trivial. An eigenvalue of order at least $p^2$
has at least $p(p-1)/2$ square-unit translates.
Otherwise the eigenvalues are $p$th roots, and Bruhat makes
$\langle U_p^aW_pv,W_pv\rangle$ constant for $a=1,\ldots,p-1$.
Fourier inversion gives equal weights at all nontrivial roots.
If these vanish, $v,W_pv$ span a $W_p$-stable Borel-fixed space.
Both generating unipotents act trivially. Thus an active factor has
at least $p-1$ active dimensions besides its neutral vector. At two
one active and one neutral dimension suffice.

Every phase occurs in the cyclic module since $sv_0$ has positive
entries. Local phases combine independently in a tensor product,
so distinct maximal sets $F$ require orthogonal summands. Each
summand witnessing $F$ is supported where $\prod_{p\in F}p\mid b_i$,
by orthogonality to the invariant coordinate lines with $p\nmid b_i$.
Its simultaneous active and neutral spaces have the asserted product
dimensions. Their unions of coordinate supports dominate the
corresponding sums of dimensions, giving the reservations by Hall's
theorem. Charge fixes the vacuum cyclic module pointwise, so in odd
dimension supports are counted in duality orbits. Finally a fixed $J$-phase space without $\ell$-activity
has scalar upper and lower $\ell$-unipotents, hence scalar Weyl.
The line criterion forbids $\ell\mid b_i$ there.
\end{proof}
Summing over sectors gives at least $C\prod_{p\in J}p$ positions
for simultaneous $J$, and at least $C\sum_{p\in P}p$ for their
union, since a product of distinct primes is at least their sum.
\paragraph{Two integral Weyl columns.}
For odd allowed primes $J$, put $F=\prod_{t\mid M,\,t\nmid q,\,t\notin J}t$
(the product is over primes) and $z=(\prod_{p\in J}W_p)v_0$.
Write $z'=W_J^{-1}sv_0$ for the complementary Weyl column.
Under Galois conjugation each local Weyl acquires a torus factor.
On $J$ the vacuum character is trivial and on the complement it
is the character of $\sqrt q$. Thus $z,\sqrt q\,z'$ are rational.
For either column $x$ the selecting CRT exponent $k$ gives
$f_i|x_i|^2=L_k^0(f_ie_i)$ by Bruhat. This is integral and is
divisible above $t$ for each complementary $t\mid f_i$, by
\eqref{eq:rotation-entry} and $(f_ie_i)^t=f_i^{t-1}(f_ie_i)$.
Since $q$ is square-free, $a_i=b_i z_i$ and $a'_i=\sqrt q\,b_i z'_i$
are integers. The complementary Weyl takes $z$ to $sv_0$ while
preserving each $J$-phase fiber $E$, whence
\begin{equation}\label{eq:phase-column}
 (b_i,F)\mid a_i,\qquad
 \sum_{i\in E}d_i^2=q\sum_{i\in E}d_i^2a_i^2 .
\end{equation}
In each fiber, divide by the squared degree gcd. If some $b_i$ is
divisible by $t\mid F$, the terms at maximal $v_t(b_i)$ sum to
zero modulo $t^2$. They number at least four for $t=2$, three
for $t\equiv3\pmod4$, and two otherwise.

\begin{lemma}[Moments on a fiber]\label{phase:moments}
Let $p$ be odd and allowed. In any complementary fiber, write $J$
for the $p$-active mass in weights $w_i=b_i^{-2}$ and $C_0$ for
the weight at one primitive $p$th root. On the neutral labels with
$p\mid b_i$ there are $a_i\in p\ZZ$ satisfying
\[
 \sum_i(1-a_i)w_i=C_0,\qquad\sum_i(a_i^2-1)w_i=J.
\]
In odd dimension these identities also count dual pairs in a fiber
excluding the vacuum. If $k$ pairs are active and $c$ have one
primitive phase, $k\equiv c\pmod2$ and $k>0$ implies
$k\ge2(p-1)$. An active three-fiber has at least eight eligible pairs.
\end{lemma}
\begin{proof}
Take $J=\{p\}$ above and use the complementary coefficients $a'_i$,
writing them as $a_i$. They vanish at active labels and are divisible
by $p$ where $p\mid b_i$. On singleton blocks $p\nmid b_i$ the
local Bruhat relation gives $a_i=1$. Set $v=sv_0$, $z=W_p^{-1}v$.
For the complementary projector $E$, the norm equality
$\|Ez\|=\|Ev\|$ gives the second identity. The vacuum Bruhat
identity for $s^{-1}Es$ identifies $\langle z,Ev\rangle$
with the Galois mean of $\sum_{i\in E}w_i\theta_{i,p}/q$.
The root means $1,-1/(p-1),0$ at orders $1,p,>p$ and equality
of primitive masses give the first identity. Both columns are
charge-even. In odd dimension $w_i\equiv1\pmod8$, giving the parity
claim. With only primitive phases their counts are equal modulo eight.
Their total is even and therefore each count is even. A higher
phase contributes $p(p-1)/2$ pairs. At $p=3$ a sole three-pair
higher orbit violates parity. At three a neutral pair contributes
$0,3$ or $7\bmod8$ to the second sum, so $4,5,6,7$ active pairs
need respectively $4,3,2,1$ additional pairs.
\end{proof}
\begin{corollary}\label{result:mixed}
If two allowed primes $p,t$ occur together, with $p$ odd, at least
$h(t)(p-1)$ active denominators are multiples of $pt$, where $h(2)=4$
and $h(t)=3$ or $2$ as $t\equiv3$ or $1\pmod4$. 
\end{corollary}
\begin{proof}
In a nontrivial $t$-phase, equal primitive masses or the orbit of a
higher phase give at least $p-1$ nontrivial $p$-phases. In each
apply \eqref{eq:phase-column} with $J=\{p\}$ at its maximal
$t$-valuation.
\end{proof}
\begin{corollary}\label{result:fiber-squares}
In square dimension simultaneous activity at three and at $p=5$ or
$7$ requires at least $18$ or $22$ labels, respectively, active at
one of these primes and with $3p\mid b_i$.
\end{corollary}
\begin{proof}
Adding twice the first fiber moment to the second expresses $J+2C_0$
as $\sum((1-a_i)/b_i)^2$ over the eligible neutral labels.
In each complementary fiber meeting $p$-activity there are at least $p-1+s_p$ eligible
labels, where $s_5=3$, $s_7=2$. Indeed a higher phase costs
$p(p-1)/2$ labels. Otherwise each primitive phase has positive equal
mass $C_0$. If each contains at least two labels the required count
already holds. A singleton gives $C_0=b^{-2}$, and the neutral
terms then express $p+1$ as a sum of rational squares. Eight is
not a square, and six is not a sum of two rational squares
(reduce a primitive integral representation modulo three).
At least two nontrivial three-phases meet $p$-activity by equal
primitive masses or the higher-phase orbit. Each gives $p-1+s_p$
labels divisible by $3p$. Conversely at least $p-1$ nontrivial
$p$-phases meet three-activity, and the three-fiber norm in each
requires a further three-neutral label with $3\mid b_i$. This gives
$2(p-1+s_p)+(p-1)$. At equality for $p=7$ the two nontrivial
three-phases each contain six singleton primitive seven-phases and
two seven-neutral labels. The other six are seven-active and three-neutral.
\end{proof}
In square dimension a $p$-phase fiber of at most two
members has $|a_i|=1$ throughout by \eqref{eq:phase-column}.
A vanishing coefficient on a two-label fiber would make $a^2-1$
a nonzero rational square for an integer $|a|\ge2$. Thus without
a $p$-power codegree there are at least $3p$ labels. If $W$ denotes
the mass at each primitive phase and $H$ the mass at higher phases,
$G=\sum d_i^2\theta_{i,p}=D-pW-H$, by square-unit sums.
The vacuum Bruhat identity $|G|=D|z_0|$ gives $MF\mid G$.

\begin{lemma}[Small blocks]\label{block:defectone}
A full $p$-block containing a label with $v_p(f)=1$ has two
defect-one tame inertia orbits with equal complementary phases.
At two it is an equal-degree pair. At odd $p$, let the sizes
and degrees of the neutral and active orbits be $(a,u),(b,v)$.
Then $a,b\mid p-1$, $(p-1)/2\mid b$, and $bv^2=(p-1)au^2$.
At three the blocks have denominators $(z,z,z)$ or $(z,2z,2z)$.
At five $(a,b,v/u)$ is $(1,4,1),(2,2,2)$ or $(4,4,2)$.
If exactly four primes divide $D$, two of them with valuation two,
both these primes are isolated in the twist graph and $D$ is square.
Here primes are joined when they divide a common twist order.
\end{lemma}
\begin{proof}
At odd $p$ the defect-one symmetric order theorem and the slice
mass identity give two tame branches, one neutral and one active,
with $bv^2=(p-1)au^2$ also on every complementary fiber
\cite[Proposition~2.13, Corollary~3.10, Proposition~19.11]{BDPcore}.
Square-unit covariance gives $h_p\mid b$ and the lists follow.
At two the defect-one branch is unramified. A second branch of
degree $e=2^k$ has $v_2(f)=k+1$ and order conductor exponent one
\cite[Lemma~2.12]{BDPcore}. For $k>0$ its field is
$F(\zeta_{2^{k+1}})$ over the unramified base $F$, by conductor
divisibility and degree. Its different exponent $ke$ would give
order conductor exponent $e(k+1)-ke=e>1$. Thus both branches are
singletons and the plane rule applies.
For the last claim call the four primes $p,q,u,v$ with $u,v$
of valuation two. A $u$-active row vanishes on all $u$-divisible
degrees. Orthogonality to the vacuum modulo $q$ supplies a degree
free of $u,q$, hence of the form $p^iv$ ($i>0$). Likewise a
degree $p^ju$ occurs. Every $p$-free nonunit in $P_q$ must meet
both, and is divisible by $uv$. A $q$-active row hits such a
nonunit by orthogonality on $P_q$ modulo $p$, so cannot be active
at $u,v$. Interchange $p,q$. A row active at both $u,v$ would hit
only nonunits divisible by both $p,q$, impossible by orthogonality.
On a nontrivial $u$-phase all complementary local groups now act
by scalars, so an active label there has $u$-power codegree,
necessarily $u^2$. Hence $D$ is square.
\end{proof}

\section{Small codegrees and isolated primes}
For a nonunit row of degree $n$, codegree $f$, write $c_Y$ for
the Galois mean of $|\chi_X(Y)|^2$. The vacuum modular relation,
row norm and lcm integrality give, respectively,
\begin{equation}\label{st}
 \sqrt f\le n^{-1}\sum_Yd_Y\sqrt{c_Y},\qquad
 \sum_Y c_Y=f,\qquad c_Y\ge(d_Y/(d_Y,n))^2\quad(c_Y>0).
\end{equation}
The nonunit hit degrees have gcd one by vacuum orthogonality.

\begin{lemma}[Small codegrees]\label{small}
No codegree is square-free or equal to four or nine. If $f=p^2$,
the nonunit degrees at entries of nonzero absolute trace are prime
to $p$, divide the $p$-free part of $n$, and have gcd one.
For $f=p^a$, $p$ odd, let $[K_X:\QQ]=p^sm$, $(p,m)=1$. Then
\[
 4\le p^{a-s-1}(1+(p-1)/m),\qquad
 \operatorname{ord}(\theta_X)\mid p^{a-1}.
\]
\end{lemma}
\begin{proof}
If $f$ is square-free, act on its coordinate vector by the local
Weyls for $\ell\mid f$ in increasing order. All complementary
factors preserve the line. At a defect-one step the degree changes
by a ratio $k^{\pm1}$ with $k^2\mid2(\ell-1)$, or $k=1$ at two,
by Lemma~\ref{block:defectone}. All remaining prime valuations are
unaffected. At a path to the vacuum, $n$ therefore divides a product
of these $k$'s. For $f>1$ its largest prime then has valuation one
in $D$, impossible by square-fullness. For $f=1$ this gives $D=1$.

At $f=4$ all hits lie in $P_2$, where nonunit degrees share a prime
in pairs. Their gcd is one, so there are at least three distinct
nonunit hit degrees. The norm budget forces exactly three hits,
each with $c_Y=1$. All three divide $n$ by \eqref{st}. If one
equals $n$, the other two would be coprime. Thus all three are
proper, and $2\le1/n+3/2<2$ by the first inequality in \eqref{st}.

At $f=p^2$ the nonzero-trace degrees have gcd one by the
trace-vector identity \cite[Proposition~5.4]{BDPcore}. Each
nonunit absolute trace $t$ has $|t|<p$. Its square costs at least
$t^2$ in the projected-row identity \cite[(2.4)]{BDPcore} besides
the vacuum. Principality gives $t\equiv0\bmod p$ when $p\mid d_Y$,
since the row field is totally ramified at $p$. Thus all the
nonzero-trace degrees divide the $p$-free part of $n$.

For $f=p^a$, average the row over wild inertia to get $v$ and
over the whole Galois group to get $w$. Their squared norms are
$p^{a-s},p^{a-s}/m$ by orthogonality. Wild averaging a $p$-power
root fixes it or kills it, so each entry $\beta$ of $v$ is
$\sum_{k=1}^{p-1}c_k\zeta_p^k$ with integral coefficients. Then
\[
 \bigl(C(\beta)+(p-1)|\langle\beta\rangle|^2\bigr)/p
       =\sum_k c_k^2/(p-1).
\]
Multiplied by the column orbit size, the right side is integral.
The $c_k$ are constant on stabilizer cosets, whose index divides
that size. At least four orbits contribute positively, since the
nonunit degrees of nonzero mean have gcd one, lie in $P_p$, and
share a prime in pairs. Summation gives the claimed inequality. It excludes
$K_X\supseteq\QQ(\zeta_{p^a})^+$, which would give $s=a-1$ and
$m\ge(p-1)/2$. By square covariance the phase order is at most
$p^{a-1}$.

A codegree-nine label would have denominator three, excluded
by Proposition~\ref{result.One-packet} below using only the prime-square
estimates just proved.
\end{proof}

\subsection{Weyl columns in square dimension}
Assume $q=1$. A nontrivial phase on a twist component has only
vacuum phases outside that component. The complementary upper and
conjugate lower unipotents act by their vacuum scalars on this
fiber. Their Weyls act by scalars, so the codegree there is
supported on the same component.

\begin{lemma}[Component columns]\label{result:component-column}
Let $J$ be a nonempty union of odd twist components, $W=W_J$, $w=Wv_0$.
Write $M_J$ and $u_i$ for the $J$-parts of $M$ and $b_i$,
respectively. Then $\epsilon=\gamma^3=\pm1$ and $g=1-w_0$ obeys
\[
 g=\sum_{\theta_{i,J}\ne1}b_i^{-2}
     \left(1-\frac{\mu(m_i)}{\varphi(m_i)}\right)>0,\qquad
 (M/M_J)(M_Jg)\equiv-\epsilon\pmod {M_J},
 \quad m_i=\operatorname{ord}(\theta_{i,J}).
\]
Its reduced denominator is $M_J$. At $i\ne0$,
$\epsilon w_i=A_i/u_i$ with $|A_i|<u_i$ and
$A_i(b_i/u_i)\equiv1\pmod{u_i}$. When $b_i$ is supported on $J$,
$\epsilon w_i=1/b_i$. Moreover
\begin{equation}\label{eq:component-norm}
 \sum_{i\ne0}w_i^2=2g-g^2,\qquad
 \sum_{i\ne0}w_i/b_i=g(\epsilon+1/M).
\end{equation}
For $J=\{p\}$ the principal $p$-block is $\{i:p\mid b_i\}$.
\end{lemma}
\begin{proof}
Write $V$ for the complementary Weyl. The split-column identities
\cite[Theorem~6.1, Lemma~6.2, Proposition~6.3]{BDPcore} use only
\eqref{eq:rotation-entry}, congruence covariance and Bruhat. In
particular $s=V+\epsilon(W-I)$, $W^{-1}v_0=w$, and the rational
columns $w,Vv_0$ are integral outside their respective prime sets.
On an active coordinate $Vv_0=0$, so $\epsilon=\pm1$.
The vacuum Bruhat sum, averaged over units, gives the displayed
Ramanujan formula. Comparing the split column with
$(sv_0)_i=1/b_i$ gives the congruences, including the exact
denominator at the vacuum (also $b_iw_i\in\ZZ$ by
\eqref{eq:phase-column}). Thus $w_0\ne0$ and $|w_i|<1$. A pure
coordinate line is invariant under $V$, so $|w_i|=1/b_i$ and
the congruence fixes its sign. The nonzero entries and the
principal height bound now give the singleton block assertion.
Finally $\|w\|=1$ and $\langle(W-I)v_0,sv_0\rangle=\epsilon g$
give \eqref{eq:component-norm}.
\end{proof}

\begin{corollary}\label{result:mixed-neutral-rank}
With at least two twist components, at least three distinct nonunit
degrees occur on twist-one labels whose denominators meet two
components. In odd dimension these account for at least six labels.
\end{corollary}
\begin{proof}
Apply Lemma~\ref{block:witness} to the partition by components.
An active label has its entire codegree supported on its own component.
In odd dimension each of the three degrees occurs on a dual pair.
\end{proof}

\begin{lemma}[Cost at an isolated odd prime]\label{result:local-entry-mass}
At an isolated odd prime $p$ let $E\le E_0$ be the reciprocal-square mass of the active labels,
and $b_i=u_i k_i$ a set of mixed nonunit denominators with
nontrivial $p$-parts $u_i$. For $x=\min\{E_0,(p^2-1)/(2p^2)\}$
one has
\[
 C'=\sum_i\frac{\|k_i^{-1}\|_{u_i}^2}{u_i^2}\le
 \frac{p+1}{p-1}x-\frac{p^2}{(p-1)^2}x^2 ,
\]
where the double bars denote the least absolute residue. For $h$
active denominators at least $B$ this gives
\begin{equation}\label{eq:cost-large-valuation}
 \sum_{i\ne0:\ 1<u_i<b_i}u_i^{-2}<
             \frac{p+1}{p-1}\frac{h}{B^2}.
\end{equation}
\end{lemma}
\begin{proof}
The active entries have pure denominator. Thus
$C'+E\le2g-g^2$, $0<g\le pE/(p-1)$ for $g=g_p$ by
Lemma~\ref{result:component-column}. Maximize
$(2-(p-1)/p)g-g^2$ on this interval for the first bound,
and drop its negative quadratic term for the second.
\end{proof}

\begin{lemma}[Heights from active levels]\label{result:activity-height-gap}
Let $p^a$ be the largest active denominator, $p^b$ the preceding
one ($b=0$ if absent), $H=v_p(M)$ and $h_a$ the activity at $p^a$.
Put $\delta=2$ in odd dimension and $1$ otherwise. Then $H\le2a-1$
and
\[
 h_a\ge\begin{cases}
 \delta(p-1)p^{2a-1-H},&H\ge2b,\\
 \delta(p-1)p^{2(a-b)}/2,&H<2b.
 \end{cases}
\]
In particular
\begin{equation}\label{eq:sparse-height}
 h_a<\delta p(p-1)\ \Longrightarrow\ H=2a-1 .
\end{equation}
If the total activity $h<\delta p(p-1)$ and $B$ is its smallest
denominator, then each nontrivial $p$-part of a nonunit denominator
is at least $B$.
\end{lemma}
\begin{proof}
If the two primitive multiplicities at $p^a$ are $\delta u,\delta v$,
equality of primitive masses gives $p^{2(a-b)}\mid u-v$.
An imbalance already costs $\delta(p-1)p^{2(a-b)}/2$ labels. At
level $p^j$ write its contribution to $g$ as $pc_j/p^{2j}$,
$c_j=h_{j,1}/(p-1)+h_{j,>}/p$, with counts at order $p$ and
above. Then $c_j$ is $p$-integral by the square-unit orbits. In a
balanced class $c_j$ is a positive multiple of $\delta$, and costs
at least $(p-1)c_j$ labels. Since $g$ has denominator $p^H$, and
lower terms have denominator at most $p^{2b-1}$ (or are absent),
$H\le2a-1$ and $v_p(c_a)=2a-1-H$ when $H\ge2b$,
$ v_p(c_a)\ge2(a-b)$ otherwise. This gives the two bounds.

Under the total-activity hypothesis the same descending argument
balances every primitive level. There is at most one minimal
order-$p^2$ orbit (counting $\delta$ copies), and no higher phase.
Write $h=\delta((p-1)j+p(p-1)k/2)$ with $k=0$ or $1$.
$j\le p-1$ or $(p-1)/2$, respectively. Thus
$2g-E\le\delta((p+1)j+p(p-1)k/2)/B^2
\le\delta(p^2-1)/B^2$. An inactive denominator with $p$-part
$1<u<B$ costs at least $\delta/u^2\ge\delta p^2/B^2$ after
subtracting active entries from \eqref{eq:component-norm}, impossible.
\end{proof}

\begin{lemma}[Valuation one]\label{result:weyl}
Suppose $p\Vert M$ is odd and isolated, and set $n=M/p$. The
$p$-active labels are of degree $n$. Each primitive phase has the
same multiplicity $0<j<p$. Off the vacuum set $a_i=p\epsilon(w_p)_i$.
Then
\begin{equation}
 a_i\equiv d_i/n,\quad nj\equiv-\epsilon\pmod p,\qquad
 \sum a_i^2=2pj-j^2,\quad \sum (d_i/n)a_i=pj+\epsilon j/n .
\end{equation}
The nonzero coefficients occur exactly at $p$-free degrees. Their
proper nonunit degrees have gcd one and share a prime in pairs.
\end{lemma}
\begin{proof}
Isolation gives active codegree $p^2$, hence order $p$ by the tame
row bound in Lemma~\ref{small}. The primitive masses agree,
so $g=j/p$ and the component lemma gives the congruences and sums.
The trace support of a prime-square row has gcd one and consists
of divisors of $n$, so its proper degrees still have gcd one and
are pairwise noncoprime by principality. The active mass gives
$j\le p+1$. At $j=p+1$ the $(p-1)j$ active unit coefficients
exhaust $2pj-j^2$, leaving no proper hit. Finally $p\nmid j$ by
the denominator of $g$.
\end{proof}

\begin{proposition}[A prime denominator]\label{result.One-packet}
If a prime denominator $p$ occurs then $p\ge5$, $p\Vert M$ and
$p$ is isolated. In the notation of Lemma~\ref{result:weyl}, $j$ is
a proper divisor of $p-1$ and exactly $pj$ labels have degree $n$.
Their $j$ neutral members form a single Galois orbit. At every
proper nonunit $p$-free degree let $k_Y=n/d_Y$. Then
\[
 \sum a_Y^2=j(p-j),\quad \sum a_Y/k_Y=\epsilon j/n,\quad
 k_Ya_Y\equiv1\pmod p,\quad \lcm_Y k_Y=n .
\]
For
each fixed primitive phase, with label set $A$, and each neutral nonunit $Y$,
\begin{equation}\label{eq:packet-trace}
 \sum_{X\in A}\chi_X(Y)=-\epsilon a_Y.
\end{equation}
For a neutral $p$-inertia orbit of length $h\le(p-1)/2$ one has
$j\mid h a_Y$.
\end{proposition}
\begin{proof}
Codegree four is excluded, so $p$ is odd. The height bound
(Lemma~\ref{result:height}) gives $p\Vert M$. Every active label
has $p$-free degree, whence $p$-primary twist by the prime-power
row rule. Thus Lemma~\ref{result:weyl} applies, and the local level
is $p$ since all active phases have order $p$.
All $m$ coefficients of degree $n$ equal one by the pure-denominator
identity of Lemma~\ref{result:component-column}.
Put $H=j(p-j)$ and $z=pj-m$. After removing this class,
\begin{equation}\label{eq:weyl-cost}
 \sum a_Y^2=H+z,\qquad \sum a_Y/k_Y=z+\epsilon j/n,qquad
 -\frac{H+z}{p-1}\le z+\epsilon j/n\le\frac{H+z}{p+1}.
\end{equation}
The bounds use $k_Ya_Y\equiv1\bmod p$ and $k_Y>1$.

Write $Q$ for the neutral projector and $u=\sum_{X\in A}v_X$.
The phase projection of $W_pv_0$ is $\epsilon u/p$. Thus $u$
lies in the Steinberg summand of the local vacuum cyclic module,
a quotient of the permutation module on $\mathbb P^1(\mathbb F_p)$,
whose nonconstant part is irreducible by double transitivity. The Steinberg neutral line
is spanned by $g=Q(W_p-I)v_0$, with $g_0=-j/p$. The complementary
Weyl is scalar $\epsilon$ on $u$, so $Qsu$ lies on this line.
Its vacuum entry $j/p$ gives $Qsu=-g$, proving \eqref{eq:packet-trace}.

For an inertia orbit $O$ of neutral labels of degree $d$ and
size $h\mid p-1$, put $b_X=\sum_{Y\in O}\chi_X(Y)$. These are
integers since the active row fields lie in $\QQ(\zeta_p)$, and
they are all $hd\bmod p$ by principality. In each packet their
sum is $-\epsilon h a_Y$. If $j\nmid h a_Y$, at least $j-1$
pairs differ by at least $p$, so the squared sum per packet is at
least $((h a_Y)^2+(j-1)p^2)/j$. Over all packets unitarity leaves
at most $hp^2-h^2(d/n)^2$. Thus $j\mid h a_Y$ for $h\le(p-1)/2$.
If $d=n$, $a_Y=1$, $h=p-1$, the opposite assumption would give
$(p-1)^2+(j-1)p^2\le j(p^2-p+1)$, hence $j\le2$, but both possible values divide
$p-1$. Consequently every such $h$ is divisible by $j$, and $j\mid m$.

By \eqref{eq:weyl-cost}, $|z|<2j$. At $z=\eta j$, $\eta=\pm1$,
its signed bound forces $\epsilon=-\eta$ and $jn\le p+\eta$,
hence equality by $jn\equiv\eta\bmod p$ and $n>2$.
Equality in the bound gives $a_Y=-\epsilon j d_Y$, contrary to
\eqref{eq:packet-trace} and $|\chi_X(Y)|<d_Y$. Thus $m=pj$.
The $j$ neutral members form one inertia orbit, so $j\mid p-1$.
If $j=p-1$, a proper orbit $h<j$ costs
$ha_Y^2\ge j^2/h>j$ by divisibility, while $h=j$ alone exhausts
the budget $H=j$. At least three proper degrees occur, so $j<p-1$.
This also excludes $p=3$, where $j=1$ would leave budget two.
\end{proof}

\begin{lemma}[Neutral compression]\label{result:steinberg-compression}
For a unit Borel-fixed vector $v_0$ in a unitary
$\mathrm{SL}_2(\mathbb F_p)$-module, suppose each nontrivial unipotent
eigenspace is a line meeting the vacuum cyclic module, which is nontrivial.
If $Q$ is the neutral projector and $h=Q(W_p-I)v_0$, then
$Q(W_p-I)Q=-hh^*/(1-(W_p)_{00})$.
\end{lemma}
\begin{proof}
The vacuum Steinberg contains every active line. Its complement is
trivial. Thus the self-adjoint compression is supported on its
neutral line, where its vacuum column determines the stated formula.
\end{proof}

\begin{lemma}[Pure denominators]\label{odd:inactive-pure}
An isolated odd $p$ has at least $p$ nonunit labels of pure
$p$-power denominator, at least $2p$ in odd dimension if $p$
itself is absent. In the latter case, if fewer than $p(p-1)$ labels
are active and fewer than $2p^2(1-1/(p-1)-1/(6p))$ nonunit pure
labels are inactive, one inactive pure dual pair has $p$-divisible degree.
\end{lemma}
\begin{proof}
If $p$ occurs the preceding proposition suffices. Otherwise put
$\delta=1$ or, in odd dimension, $2$. The primitive masses or the
orbit of a higher phase give at least $\delta(p-1)$ active labels.
At equality each primitive phase occurs on one label or one pair,
all of common denominator $B\ge p^2$. Then
$g=\delta p/B^2$ and $M_p=B^2/p$. The squared Weyl norm remaining
after the active entries is at most $\delta(p+1)/B^2$, so a
nontrivial $p$-part $u<B$ cannot occur. For a mixed denominator
$uk$, its signed numerator $A$ satisfies $|A|k\ge u-1$, hence
$|w_i|/b_i\le w_i^2/(B-1)$. Subtracting the pure contributions
from both identities in \eqref{eq:component-norm}, and denoting
the inactive pure mass multiplied by $B^2$ by $E$, gives
\begin{equation}\label{eq:packet-shell}
 |\delta-E+\delta\epsilon p/M|\le
 \bigl(\delta(p+1)-E-(\delta p/B)^2\bigr)/(B-1).
\end{equation}
In particular $E/\delta\ge1-1/(p-1)-1/(6p)>1/p^2$ as $M\ge6B$.
There is an extra pure label or pair.

More generally let $V$ and $P$ be active and inactive pure masses.
When all phases are primitive and every mixed nontrivial $p$-part
is at least $p^2$, the same subtraction with $g=pV/(p-1)$ gives
\begin{equation}\label{eq:unequal-pure-shell}
 P\ge g-V-g/M-(2g-V)/(p^2-1)
  =\frac{V}{p-1}\left(1-\frac pM-\frac1{p-1}\right)>0 .
\end{equation}
These hypotheses hold for $h<p(p-1)/2$ active labels, by orbit sizes
and \eqref{eq:cost-large-valuation}. They also hold in the odd
case with $h<p(p-1)$. The height lemma balances primitive counts,
rules out higher phases, and bounds each nontrivial $p$-part below
by the least active denominator $B$. Hence $V\ge2(p-1)/B^2$,
and $M_p=B_{\max}^2/p\ge pB$. If all inactive pure labels have
$p$-free degree then each has denominator $M_p$ and the assumed
count bounds $P$ strictly below the lower bound in
\eqref{eq:unequal-pure-shell}.
\end{proof}

\section{Multiplicity and prime support}
Write $P_p$ for the principal $p$-block. We use $\langle z\rangle$ for
the arithmetic mean of the Galois conjugates of $z$. A nonzero cyclotomic
integer satisfies $\langle|z|^2\rangle\ge1$, with equality precisely for
a root of unity. For a row principal at two, if all its entries on an odd-degree
class of size $s$ have modulus one, their sum is an integer of absolute value
at most $s$ and parity $s$. It is therefore a sum of $s$ signs. If it
equals $\pm s$, every entry equals that sign.

\begin{lemma}[A two-column Gram bound]\label{lem:congruent-rows}
Let a Galois-stable family of $m$ orthogonal rows of cyclotomic integers have norm
squared $f=p^a$ and vacuum entry one, and suppose they are perpendicular to the degree
row. Suppose they vanish off a common coordinate set $F$ and are
congruent to the degree row above $p$ on $F$. Then
$m\le f-\lceil\sqrt p\rceil$ and, for $a\ge2$, $m\ne f-p$.
\end{lemma}
\begin{proof}
Some $p$-free nonunit $Y\in F$ occurs, by degree orthogonality modulo $p$.
Let $w,A$ be the sums of the entries and their squared moduli there.
Both are integers and $A\ge m$ by the norm inequality. Set $k=f-m$,
$l=f-A$. The Gram matrix of $e_0,e_Y$ after projection off the rows,
multiplied by $f$, has diagonal $k,l$ and off-diagonal $-w$.
It is positive definite because proportionality would give a norm ratio
$\sqrt{l/k}\le1$, whereas the ratio of the degree components is $d_Y>1$.
Since $w\equiv-kd_Y$, $l\equiv kd_Y^2\pmod p$, one obtains
\begin{equation}\label{eq:saturation-gram}0<p\mid kl-w^2\le k^2.\end{equation}
If $k=p$, then $l>0$, $p\mid l$ and $l\le p$, whence $l=p$, $w=0$.
This holds at every $p$-free nonunit. At every other coordinate the
row sum is divisible by $p$, contradicting degree orthogonality modulo
$p^2$, since the vacuum sum is $p^a-p$.
\end{proof}

\begin{lemma}\label{result:saturation}
At degree $n$ of prime-power codegree $f=p^a$, the preceding multiplicity
bounds apply. For odd $n$ and dyadic $f$ one has $m\le f-6$ at $f=8$
and $m\le f-8$ at $f\ge16$. For odd $n$ at $f=16$, either $\omega(n)=3$ or $r<38$ implies $m\le6$.
\end{lemma}
\begin{proof}
Use Lemma~\ref{block:admissible} for the first assertion. Suppose
$f-m\le5$ at eight or $f-m\le7$ at $f\ge16$. Every even hit costs
at least $4m\ge f$ in its column, by lcm integrality, besides the
vacuum. Thus $P_2$ is exactly the odd-degree set, all of whose degrees
divide $n$. Its mass is integral at two, so $8\mid |P_2|$, since
odd squares are one modulo eight. The proper nonunit hit degrees have
gcd one and are pairwise noncoprime. Hence at least three distinct
ones occur. In particular $f-8<|P_2|\le f$. Thus $|P_2|=f$ and
every entry is a root of unity.

Put $h=f-1-m$. Degree orthogonality for each row takes the form
$1+nt+\sum_{i=1}^h\epsilon_i a_i=0$ with signs $\epsilon_i$,
$t\equiv m\pmod2$, $3\le a_i\le n/3$ the proper degrees, $n\ge15$.
For $h=4$, $t$ is odd and mixed signs give absolute proper sum
at most $n-3$, so all four signs agree and their sum is determined.
For $h=3,5$, one has $t=0$. Two solutions differ in an even number
of positions, and four differences would leave a signed degree $-1$
(or be impossible at $h=3$). Two differences only exchange equal
degrees. Thus the signed sum $t_b$ at each proper degree $b$ of
multiplicity $s_b$ is constant across rows, with $\sum|t_b|\ge3$.
Bessel on the vacuum and class-constant lines yields
$m(1+9/h)\le m(1+\sum t_b^2/s_b)\le f$, impossible in the asserted
ranges. For $h=6$, $t$ is odd. Bessel on the degree vector restricted
to the vacuum and proper coordinates gives
$m/f\le n^{-2}(1+\sum a_i^2)\le1/225+4/9+2/25<9/16$ if at most
four of them equal $n/3$. Otherwise five equal $a=n/3$, and the
remaining degree, coprime to $a$, must be three. Its singleton entry
is a sign, forcing $a\mid1\pm3$, impossible.

For the final bound assume $m\ge7$. Again $4m>16$ gives the same
sixteen odd labels. Four odd primes would reserve at least
$2+4+6+10$ more labels, so precisely three occur under either
hypothesis. The proper labels omitting each of these primes form
three disjoint nonempty sets. On a set omitting $p$, its signed
degree sum is $-1\pmod p$. On one or two labels this fixes the
signs, since changing both negates the sum. Actual entries are
fixed as well. Two fixed columns would give $3m\le16$. With at
most eight proper positions the remaining possibility includes a
singleton and a triple. On a triple two solutions differ in exactly
two positions, exchanging equal residues after orienting them up to
sign. Thus each oriented residue class has fixed sum. At least one
is nonzero. Projection onto the vacuum, singleton and these class
lines gives $m(2+1/3)\le16$.
\end{proof}

\begin{proposition}[Rows obtained from a defect-one block]\label{result:effective}
In an equal-degree defect-one $\ell$-block $B$ of codegree $\ell p^a$
($p\ne\ell$), there is a unique $\ell$-neutral label $X$. Its character
restricted to $F=\{Y:\theta_{Y,\ell}=1\}$, extended by zero, equals
$v_B=\ell^{-1}\sum_{i\in B}\chi_i$. These rows have squared norm $p^a$,
are pairwise orthogonal, supported on $P_p$, and have character field
$K_X\subseteq\QQ(\zeta_{p^a})$. Lemma~\ref{lem:congruent-rows}
applies to them. Their vacuum Gauss identity is
\begin{equation}
\sum_Yd_Y\theta_Yv_B(Y)=\eta_B\sqrt{D/\ell}\,\theta_X^{-1},
\qquad |\eta_B|=1 .
\end{equation}
The factor is constant on dual pairs and under odd $p$-inertia.
At squared norm $p^2$, the nonunit degrees in the nonzero-mean
support of an actual or such an averaged row are $p$-free, have gcd
one and are pairwise noncoprime. At least three distinct degrees
and at least three primes other than $p$ occur.
\end{proposition}
\begin{proof}
Write $s=W_\ell W_pV$ and $c=\ell^{-1/2}\sum_Be_i$.
Lemma~\ref{block:defectone} and Lemma~\ref{block:admissible} give uniqueness
of $X$, with complementary phases constant on $B$. The projection
of $se_0$ on $B$ spans $c$. Its inverse $W_\ell$-image is neutral,
so $W_\ell^{-1}c$ spans $e_X$. The central square $W_\ell^2$ preserves
that line. The factors of $V$ preserve it too, by
Lemma~\ref{block:admissible}. Hence $X\in P_p$ and $sc$ is supported
on $P_p\cap F$. $s(c^\perp\cap\mathbb C B)$ is supported off $F$.
This gives the stated restriction and the norm by unitarity.
The Galois stabilizers of $X,B,v_B$ coincide, as $s$ separates
the block sums. Inertia at $\ell$ preserves the residue block and
fixes $X$, so the conductor bound removes $\ell$. Average the vacuum
Gauss equations over $B$. The local inverse twist sum is $1\pm i$
at two and a signed quadratic Gauss sum at odd $\ell$, each of modulus
$\sqrt\ell$. The other factors are fixed under odd $p$-inertia.
For the last assertion the trace and congruence argument of
Lemma~\ref{small} applies on $F$ verbatim. Its hit labels
belong to $P_p$, where $\chi_Y(Z)\equiv d_Z\ne0$ above $p$.
Coprime vanishing therefore gives the pairwise property. A pairwise
intersecting family of prime supports with empty intersection requires
at least three primes and three members.
\end{proof}

\begin{lemma}[The height at a prime-square codegree]\label{result:height}
An actual or averaged prime-square row of odd prime $p$ has degree
$n$ prime to $p$. In particular a denominator $p$ in square dimension,
or in an equal $\ell$-block at $D=\ell M^2$, forces $p\Vert M$.
\end{lemma}
\begin{proof}
Sum the rows of its orbit to an integer vector $v$, of squared norm
$hp^2$ and vacuum entry $h$, where $h$ is the orbit size. Only
$p$-inertia is needed. Each nonzero coordinate has $p$-primary twist
and degree dividing the prime-to-$p$ part of $n$ by the preceding
proposition, principal height and local admissibility. On the orbit
the twists are trivial or of order $p$ by Lemma~\ref{small}.
Its norm argument applies also to the averaged translates, whose mean
support has three pairwise noncoprime degrees of gcd one.
The summed Gauss equations have right side
$\eta pn\sum_{X\in O}\theta_X^{-1}$. If $p\mid n$ and the orbit is
neutral, taking absolute values gives
$ pnh\le h+(n/p)(hp^2-h^2)<pnh$, since $n/p>1$.
For an active orbit apply the $p$-Galois average weighted by the
quadratic character. It annihilates roots of order different from
$p$, and sends $\zeta_p^k$ to $(k/p)\tau/(p-1)$ for the quadratic
Gauss sum $\tau$. The inverse row twists belong to one square class.
$\eta$ and $v$ are fixed. Cancelling the common factor gives
$pnh\le (n/p)(hp^2-h^2)<pnh$.
\end{proof}

\begin{lemma}\label{result:small-effective}
At codegree $4\ell$ ($\ell$ odd prime) the multiplicity is one of
$h,2h,\ldots,7h$, $h=(\ell-1)/2$, and there is no equal $\ell$-block.
At codegree eighteen the degree is three-free, its multiplicity at most
ten, and at least four dimension primes occur. In square dimension a
prime denominator $p$ satisfies $p\ge5$, $p\Vert M$ and at least four
primes divide $M$. All $p$-active labels have denominator $p$.
\end{lemma}
\begin{proof}
An equal block at $4\ell$ would yield an averaged norm-four row.
The norm-four argument in Lemma~\ref{small} applies to its principal support and vacuum Gauss identity.
In an unequal block the degree cannot be the smaller. Its square
ratio would divide $4\ell$, giving the forbidden codegree $\ell$.
Thus all these labels are active in packets divisible by $h$. Their
blocks have mass $mn^2\ell/(\ell-1)<4\ell n^2$ since they omit the
vacuum. At eighteen the dyadic blocks are equal pairs, and the
preceding lemma gives $3\nmid n$. Each averaged row meets the neutral
member of each pair and at least three distinct proper nonunit degrees
of gcd one, so $1+m/2+3\le9$. The assertion on the number of primes
follows from Proposition~\ref{result:effective}. In square dimension
also use Lemmas~\ref{small}, \ref{small} and~\ref{result:weyl}.
\end{proof}
Thus for $D=5M^2$ with $9\mid M$, all denominator-three labels are in
unequal quinary blocks, which have as many denominator-six labels.

\begin{lemma}[Exclusion of codegree eight in small rank]
The greatest degree $x$ cannot have codegree eight in rank at most $44$.
\end{lemma}
\begin{proof}
Here $D=8x^2=2M^2$. Its row meets all odd labels, with degrees dividing
$x$. Their number is divisible by four. If $x$ is odd this is the
whole support, with at least three distinct proper divisors, whence
$x\ge105$. There are eight entries, all roots of unity. There are
$m\le2$ maximal labels by Lemma~\ref{result:saturation}. At $m=1$ the
vacuum Gauss bound would give $\sqrt8\le1/x+1+4/3+1/5+1/7$.
At $m=2$ each maximal row field lies in $\QQ(\zeta_8)$ and has
at most two conjugates, so its root entries belong to $\{\pm1,\pm i\}$.
The rows are rational or complex conjugate, hence agree or oppose at
each hit. In each row the two maximal entries have even integral sum,
zero since all other degrees total at most $1+5x/3<2x$.
By symmetry both maximal hits belong to the half-difference.
Row orthogonality gives four hits in each half-row. Degree orthogonality
on the half-difference makes its two proper degrees equal to some $y$.
The half-sum has three proper degrees with signed sum $-1$. A repeated
degree $a$ would force the third to be $1$ or $2a\pm1$, contrary to
pairwise noncoprimality. The two Gauss half-sums are therefore bounded
by $1+(1/3+1/5+1/7)x<7x/10$ and $2x+2y\le8x/3$. Their squared
moduli must sum to $8x^2$, a contradiction.

For even $x$, the norm-eight calculation of
\cite[proof of Theorem~20.7]{BDPcore} uses precisely principality,
$(x/d_Y)\chi_X(Y)\equiv0$ above two and \eqref{st}. It gives
$\chi_X=A+B\sqrt2$, where outside the vacuum $A$ consists of
three signs at the odd nonunits, and $B$ of two opposite signs at
one degree $y$. The odd degrees are distinct, divide $x/2$, have
gcd one and share primes in pairs, so at least three odd primes
divide $x$. In particular $y$ is even. If $y<x$ then
$\sqrt8\,x\le1+(1+1/2+1/3)x/2+\sqrt2\,x$, impossible.
The signs at the odd labels are unique by degree orthogonality
(a change of signs would have to cancel two of the three distinct
odd degrees). Bessel gives $4m\le8$. The conjugate rows give
$m=2$. These two rows and the four odd labels have purely dyadic twist.

Odd-prime activities require $2(p-1)$ disjoint positions per prime
besides the two maximal and four odd labels. Thus the odd primes are $3<l<t$.
Four primes cost at least $44$ places, and $(5,7,11)$ costs $40$.
Each odd nonunit omits a different prime by gcd one and pairwise
noncoprimality, giving odd denominator floors $12,4l,4t,12lt$.

Write $2k$ for the multiplicity of denominator three. If $k>0$ then
$3\Vert M$. The $k$ dyadic pairs yield orthogonal effective norm-nine
rows, all meeting exactly the dyadic-neutral three-free labels, of
degrees dividing $n=M/3$. At $0$ and the unique three-free odd
nonunit $U$ their entries are $1,-A_U$ by orthogonality to $A$,
the mean of the maximal rows. Their Galois-invariant sum is integral,
congruent to $k d_Y$ modulo three at every hit $Y$, and has squared
norm $k(9-2k)$ off $0,U$. Thus $k\le4$. At $k=4$ the four neutral
$n$-labels exhaust this norm, yet any further hit would be nonzero
modulo three. The nonunit hit degrees would then have common divisor
$d_U\mid n$. At $k=3$ every entry is divisible by three, so the
sole remaining hit has degree $|1-A_Ud_U|$ by degree orthogonality,
coprime to $d_U$ in the principal three-block. Thus $k\le2$.

Reserve $2k$ threes and $4-2k$ sixes, and for $s=l,t$ reserve
$2(s-1)$ floors $s$, doubled if denominator $s$ is absent. Fill
the remaining places to $44$ at floor four. The mass bound decreases
with each prime and increases with $k$. For $(l,t)\ne(5,7)$ it is
at most $11977/6600$ (take $5,11$, $k=2$). For $(5,7)$ with at
most one of $3,5,7$ occurring as denominators, its maximum is
$3681/1960$ at $k=0$ and $57959/29400$ at $k>0$. All are below two.

Thus two primes $a,b$ in $\{3,5,7\}$ occur as denominators and
have valuation one in $M$. Writing $t$ for the third, the odd
degrees are $ab,at^i,bt^j$, with $|at^i\pm bt^j|=ab\pm1$.
For $(a,b)=(3,5),(3,7),(5,7)$ the $t$-divisible right side is
$14,20,36$, respectively. If $i\ne j$, dividing by the full $t$-part
leaves $2,4,4$, respectively, each less than $at-b$. Thus $i=j=1$. The signs are $-,-,+$ at $15,21,35$.
Choose $p\in\{5,7\}$ among the denominators. Every degree in $P_p$
is $p$-free by height, and their nonunit gcd is one. Some $Y\in P_p$
therefore omits both five and seven. Its prime support is $\{2,3\}$.
The singleton odd rows are rational. By symmetry, integrality and
coprime vanishing, $\chi_Y(15)=5u$, $\chi_Y(21)=7v$, $\chi_Y(35)=0$,
where $u,v\in\ZZ$ and $|u|,|v|<3$. Orthogonality to the maximal
mean gives $5u+7v=1$, impossible modulo seven.
\end{proof}

\subsection{Prime denominators in square dimension}
Suppose $D=M^2$.
\begin{lemma}
A denominator-four row of even degree or multiplicity at least four
hits only $4\mid b_Y$. Its nonunit hits include a multiple of each
full prime power in $M$. With three support primes a denominator
$p^a$ forces $p^2\mid M$, and four fours force $8\mid M$.
\end{lemma}
\begin{proof}
The dyadic column estimate and the height bounds are
\cite[\S21.2 and Lemma~21.6]{BDPcore}. The common-divisor assertion
also follows from gcd one on the hit degrees.
\end{proof}

\begin{lemma}\label{result:neutral-inertia}
If $v_p(D)>0$ is even, $p$ odd, and all $p$-phases have order at most
$p$, the number of active labels is at least $(p-1)(1+c/2)$, where
$c$ is the number of neutral labels less their number of inertia orbits.
For a prime denominator write $j$ for its multiplicity per primitive
phase. At five, $j=1$ and the four proper Weyl coefficients are signs.
At seven, $j=1$ or $2$. At $j=2$ the proper coefficients have magnitudes
$1,1,2,2$, the first two on a moved pair and the last two on distinct
singleton degree classes. Neither prime denominator occurs in odd
dimension.
\end{lemma}
\begin{proof}
For $G=\mathrm{SL}_2(\mathbb F_p)$, its unipotent $U$ and its Borel
$B$, $c=\dim V^U-\dim V^B$ by the permutation action of the torus.
The vacuum Steinberg supplies $p-1$ active dimensions. The other
nontrivial torus types in $V^U$ occur by reciprocity in
$\operatorname{Ind}_B^G\alpha$. A non-self-inverse character gives
two torus lines and $p-1$ active dimensions, and each of the two
quadratic constituents gives one line and $(p-1)/2$ active dimensions
\cite[Prop.~2.14]{PrasadSL}, \cite[Lem.~4.3]{Nevins}.
Indeed the induced module restricts to $U$ as
$\mathbf1\oplus\mathrm{Reg}_U$. Summing proves the bound.

The neutral degree-$M/p$ orbit uses $j-1$ of the at most $2(j-1)$
units of deficit. Each proper orbit of length $h$ and coefficient
$a$ costs $ha^2$ out of $j(p-j)$, with $j\mid ha$
(Proposition~\ref{result.One-packet}). At least three degrees occur.
At $(p,j)=(5,2)$ at most one pair moves, and the norm six pays at
most one other even coefficient. Thus $j=1$, giving four signs.
At $(7,3)$ a pair costs at least eighteen, and a triple and two
fixed labels cost at least twenty-one, against norm twelve. At
$(7,2)$ norm ten on at least three degrees forces the indicated
pair and two even singletons. The possible $j$ are the proper
divisors of $p-1$ by Proposition~\ref{result.One-packet}. Odd duality
excludes a singleton neutral row or a singleton proper degree.
\end{proof}

\begin{lemma}[The least possible activity at an isolated prime]\label{result:packet-idempotent}
Suppose an isolated odd prime $p$ has $l(p-1)$ active labels of
common denominator $B$, with $l=1$ or, in odd dimension, $l=2$.
If $B\ne p$, suppose fewer than $lp^2(1-1/(p-1)-1/(6p))$ inactive
nonunits have pure $p$-power denominator. Then $B$ has multiplicity
$lp$. If $B=p$ and $l=1$, all its rows have value $-\epsilon a_Y$
on each proper $p$-free label $Y$.
\end{lemma}
\begin{proof}
If $B\ne p$, no smaller pure denominator fits in the remaining norm
$l(p+1)/B^2$. By \eqref{eq:packet-shell} the scaled inactive pure mass
$E/l$ lies between $1-1/(p-1)-1/(6p)$ and $1+1/(p-1)+1/(6p)<2$.
Each label away from $B$ costs at most $1/p^2$. Hence $B$ has an
inactive label. It has exactly $l$, by the upper bound and duality.
At $B=p$ the multiplicity is Proposition~\ref{result.One-packet}.
Assume $l=1$. There is one neutral
member $X$. The active rows have the asserted value by
\eqref{eq:packet-trace}. The vacuum Steinberg already uses all
$p-1$ active dimensions, so its complement is trivial: there $U_p$
is the identity, as are all its conjugates. Let $Q$ be the neutral
projector. The self-adjoint $Q(W_p-I)Q$ has rank one, since the
Steinberg neutral space is a line. For $w=W_pe_0$ its vacuum
entry is $w_0-1=-1/p$ and $w_X=\epsilon/p$ by the local column
formula. Hence $(W_p)_{YX}=w_Yw_X/(w_0-1)=-\epsilon w_Y$.
Local admissibility and Bruhat make the complementary Weyl act by $\epsilon$ on $X$, giving
$\chi_X(Y)=p s_{YX}=-p w_Y=-\epsilon a_Y$.
\end{proof}

\begin{lemma}\label{result:cross-weyl}
At a prime denominator $p$, let $a_Y$ be its proper coefficients.
For any nonunit row $X$ with $p\nmid b_X$,
$\sum a_Y\chi_X(Y)=\epsilon j$. If $m$ such rows of equal codegree
$f$ have all their proper hits in $Z$, then
$\sum_Z a_Y^2\ge mj^2/(f-m)$. With twelve denominator-four rows
and a denominator seven, or eleven such rows and exactly seven sevens,
there are six proper seven-free labels, with denominators $7k$,
$k\equiv\pm1\pmod7$, $k\ge6$.
\end{lemma}
\begin{proof}
At coprime denominators the nonunit entry of $s$ vanishes by lcm
integrality and strict unitary modulus at every embedding.
$W_pe_X=e_X$ by admissibility and Bruhat. Expanding
$\langle W_pe_0,se_X\rangle=1/b_X$ gives the identity.
Project $\sum_Za_Ye_Y$ and $\epsilon j e_0$ off the $m$ normalized
rows and apply Cauchy--Schwarz. Equality forces
$\sum_Za_Yd_Y=-\epsilon jm/(f-m)$ by taking the degree component.
Thus at twelve fours all four proper five-signs satisfy $4\mid b_Y$.
Three would be equality and the omitted signed degree would be
$\pm4$, impossible.

At twelve fours mass forces $j=1$ at seven. The norm six on at
least three degrees permits six signs or $2,1,1$ in magnitude.
In the latter case all three are singletons, hence their columns
are rational. Exactly one sign column $A$ has odd degree by the
signed degree sum. Call the other $B$. The four-row entries at $A$
are odd integers and the column norm forces signs. Their sum $w$
satisfies $|w|<16-m$ by \eqref{eq:saturation-gram}, so both signs
occur at least $m-7\ge4$ times. Column $B$ is even and at most
three entries can be nonzero. The double coefficient times its
column has fixed residue modulo four by dyadic principality, so
the cross identity modulo four forces a $B$-entry for every occurrence
of one sign at $A$, a contradiction. The congruence $ka_Y\equiv1$
gives the asserted indices.
\end{proof}

\begin{lemma}
If $M$ has prime support $2,t,5,7$, five and seven occur as denominators,
seven has one packet, and at least eight fours occur, their proper
supports have respectively four and six signs and meet in a singleton.
\end{lemma}
\begin{proof}
The local columns give $\sum_{H_5\cap H_7}a_{5,Y}a_{7,Y}=-1$, where
$H_p$ is the proper support. Two five-free witnesses omit two and $t$
and are outside $H_7$, as nonunit degrees are no prime powers. Hence
at most two labels are shared. With signs, the identity gives one.
If the seven-coefficients are $2,1,1$ in magnitude, their singleton
columns omit $2,t,5$, respectively in some order, and the shared
column $C$ has magnitude one. The double $Y$ has even degree, since
eight odd integer entries together with its seven-packet mass $4/7$
would exceed unit column norm. Let $A$ be the odd sign column.
There is a unique odd five-free witness $U$ by parity of its signed
degree sum and the two distinct even witnesses omitting $7,t$.

Multiply each of these four columns $V$ by $\epsilon\operatorname{sgn}a_{p,V}$
at its respective prime ($C$ at seven). By Lemma~\ref{result:packet-idempotent},
on columns $0,A,Y,U$ the seven- and five-packets give the rows
$v=(1,-1,-2,0)$ and $w=(1,0,0,-1)$.
On eight four-rows the $A,U$ entries are signs, and the even $Y$
column has at most one nonzero entry, of magnitude two, by the column
norm. Write $r_i=(1,a_i,y_i,u_i)$ for these rows, with $y_2=\cdots=y_8=0$.
Removing them and the packets gives
\[
 G=560I-80v^tv-112w^tw-35\sum_{i=1}^8r_i^tr_i\succeq0.
\]
There are $12\binom{10}3=1440$ choices of $r_1$ and the multiset of
the seven pairs $(a_i,u_i)$, $i>1$. Testing the fifteen principal
minors leaves only $r_1=(1,-1,2,-1)$ and multiplicities $(0,1,2,4)$
in lexicographic sign order. Then $G(6,-1,-7,1)^t=0$. Pairing with the vacuum row gives
$6-\tilde d_A-7\tilde d_Y+\tilde d_U=0$ for the signed degrees,
and the seven-column identity gives $\tilde d_A+2\tilde d_Y+\tilde d_C=1$.
Elimination contradicts $t\mid\tilde d_A,\tilde d_C,\tilde d_U$ and
$t\ne5$.
\end{proof}

\begin{lemma}\label{result:small-Gauss}
Write $T_b=\sum_{b_i=b}\theta_i$. For complete classes indexed by
$J$, of denominator lcm $L$, let $R$ be the complementary reciprocal
mass after deleting classes with $T_b=0$. Then
$L^2(R+1/M)<1$ implies $\sum_{b\in J}T_b/b^2=0$.
With denominator five present, eleven fours permit at most three
sixes. Twelve fours permit no six and at most two eights. In the
latter case the four proper five-free denominators are multiples of
twenty at least eighty. Two eights then give $T_8=0$ and preclude a
singleton nine.
\end{lemma}
\begin{proof}
Every conjugate of the integral sum $\sum(L/b)^2T_b$ is bounded by
$L^2(R+1/M)$ in the vacuum Gauss identity. Now $T_5=0$, $M\ge420$,
and eleven fours and four sixes leave $1/720$, forcing
$9T_4+4T_6=0$ contrary to reduction at two. With twelve fours the
remaining mass is $1/20$, giving $T_4=0$. A six would require
$1/M\ge2/36-1/20$. The indices $k_i$ of the five-signs are
multiples of four, congruent to the signs modulo five. Below sixteen
only $k_i=4$ with sign minus is possible, and even three terms of
size $1/16$ cannot offset $-1/4$ to $\pm1/n$. Three eights
would give $T_8=0$ with odd cardinality. With two eights the house
of $T_8$ is at most $64(3/160+1/420)<\sqrt2$, forcing zero by the
dyadic mean bound. A singleton nine then requires
$2/81\le3/160+1/M$.
\end{proof}
When $m_4$ is even, Lemmas~\ref{result:coupled-gap}
and~\ref{result:coupled-gap} give $T_4=T_6=0$ for
$m_6\le3,\ 36(R+1/M)<1$, or $m_6\le8,\ 72(R+1/M)<1$.
They give $T_4=0$ for $m_6\le5,\ m_8\in\{0,2\},\ 96(R+1/M)<1$.
Here $R$ omits exactly the classes under consideration together with
known zero sums.

\begin{lemma}\label{result:joint-projection}
Select $s_q$ rows of codegree $f_q$ a power of $q$, and let $p$ be
another prime. Take an orbit union $U$ of $p$-free nonunits containing
all such columns of nonzero Galois mean on these rows. Choose positive
orbit-constant $w_Y$ prime to $p$, with $q\mid w_Yd_Y$ for each $q$.
For $1\le a<p$ put $c_Y=w_Y[ad_Y/w_Y]_p$ on $U$ and zero elsewhere,
where brackets denote least absolute residues. Choose
$0<u_q<pq$ with $q\mid u_q$, $u_q\equiv-a\pmod p$. Then
\[
 \|c\|^2\ge\sum_q(s_q/f_q)\operatorname{dist}(t+u_q,pq\ZZ)^2-t^2
 \quad(t\in\mathbb R).
\]
For one class set $l=u-pq$, $h=\max(1,(p-1)/2)$,
$S_j=(s-j)l+ju$ and $Q_j=(s-j)l^2+ju^2$. Then
\[
 \sum_Uw_Y^2\ge H_f(p;q,s):=
 \left\lceil\frac{\sum_{a=1}^h
 \left\lceil f^{-1}\min_{0\le j\le s}(Q_j+S_j^2/(f-s))\right\rceil}
 {\sum_{a=1}^h a^2}\right\rceil .
\]
The minimum is taken over integers $j$. When rows and columns are dual pairs,
take $j$ even and round both ceilings upwards to even integers.
\end{lemma}
\begin{proof}
The row evaluations $z_X=\sum c_Y\chi_X(Y)$ are integers by Galois
symmetry, congruent to $-a\pmod p$ by degree orthogonality and to
zero above $q$ by principality (at $p=2$ take the odd residue to be
one). Apply Bessel to $c+te_0$. Optimizing $t$ on one class gives
$f\|c\|^2\ge\sum z_X^2+(\sum z_X)^2/(f-s)$.
At fixed sum the least square cost uses neighboring residue levels.
If all levels are on one side of zero, moving them toward zero
decreases both terms. This gives the two lifts $l,u$. Summing over
$a$ permutes the absolute residues at each coordinate. Duality
pairs the real evaluations and the coordinate costs.
\end{proof}
When $U$ contains \emph{all} $p$-free nonunits and $w_Y=1$, every
row, including the vacuum, has $z_X\equiv-a\pmod p$ by the dimension
identity. Their distribution with weights $1/f_X$ has mean zero and
second moment $\|c\|^2$. Let $L$ interpolate the square linearly on
consecutive elements of $-a+p\ZZ$. If the selected rows have
evaluations in $q\ZZ$, one obtains
\begin{equation}\label{eq:full-census}
 \|c\|^2\ge f^{-1}\min_{0\le j\le s}
       \{Q_j+(f-s)L(-S_j/(f-s))\}.
\end{equation}
Indeed the residual rows have prescribed mean. Outside the range
of sums obtained with $l,u$, moving a level towards zero lowers the
total cost: take $a\le(p-1)/2$, as suffices by sign symmetry. Below
$s l$ the residual mean is positive where $L$ is increasing. Above
$s u$ a step upwards costs at least $pq(2u+pq)$, exceeding any
saving at negative mean, since the upper slope there is $p-2a$.
In the general one-class application with a common row-sum parity
one may minimize over both parities. For $m$ four-rows of even
degree, weighting their nonunit $p$-free hits by $1$ at even
degrees and $2$ at odd gives $e_p+4o_p\ge H_{16}(p;2,m)$.

\begin{lemma}
In odd dimension a five- or seven-power row forces at least six
nonunit labels of degree prime to its prime. For seven the lower
bound is eight if the union of the other principal blocks has mass
above $3/7$. With exactly four eleven-free nonunits and a three-power
row, one of those degrees is three-free if $11\mid D$.
\end{lemma}
\begin{proof}
At any dimension prime
$p$, averaging $(z_X+a)(z_X+a-p)\ge0$ gives
$\|c\|^2\ge a(p-a)$. Equality confines every $z_X$ to
$\{-a,p-a\}$. A row of five- or seven-power codegree forces at least
six nonunit labels of degree prime to its prime. Indeed the dimension
modulo $p$ excludes fewer than two dual pairs. At two pairs the
absolute residues are $(1,1)$ at five, $(1,3)$ at seven. Equality
holds with $a=1,2$, respectively. The even augmentation has value
$4,-2$, respectively. Its half yields coprime degrees (odd degrees
with signed sum two at five, or a weighted sum $\pm1$ at seven),
contrary to principality. At eleven two pairs have residues $1,2$:
equality at $a=1$ yields values $-1,10$, so a three-power row implies
that one pair is three-free. At seven three pairs have residues
$(1,1,1)$ or $(2,2,3)$. Equality holds with $a=1,3$ respectively.
The even augmentation $7-a$ has mass $a/7$, since the mean is zero.
Every row in another principal block takes that same value by reduction
at its prime. Thus, if a seven-power row occurs and the union of the other
principal blocks has mass exceeding $3/7$, at least eight nonunit labels
are seven-free.
\end{proof}

\subsection{The support of an even square dimension}
In this subsection $r\le44$, $M$ is even, and $m$ counts denominator-four
labels. A reservation $(h,b)$ charges $h$ distinct positions of denominator
at least $b$. After disjoint reservations $(h_i,b_i)$, with floor $c$
elsewhere, the upper mass is
\[
 U((h_i,b_i);c)=\sum h_i/b_i^2+(44-\sum h_i)/c^2.
\]
A negative remaining count or an upper mass below one excludes the
profile. Padding by unreserved positions covers every smaller rank.

\begin{lemma}[Three-prime support]\label{result.Three-prime-rank}
At least four primes divide $M$.
\end{lemma}
\begin{proof}
Suppose the primes of $M$ are $2<p<t$. Each occurs to
exponent at least two. A $u$-free witness in $P_v$ ($u\ne v$)
has both other primes in its degree and $v$ in its denominator.
Thus $M\ge900$ and no prime denominator occurs. Write $m\le13$ for
the multiplicity of four. Other nonunit denominators have floor six
(eight if $p>3$). We use the following position counts to bound the
sum of reciprocal squares, padding to $44$ places.

Put $A=4$ if $v_2(M)=2$, and $A=8$ otherwise. The first case has
$m\le3$. Odd degrees occupy $3\bmod4$ nonunit places. These include
the $m$ fours when $A=4$, and one further denominator divisible by
$At$ from $P_t$. The other odd places have floor $4p$ if $A=4$,
and $8$ otherwise. For $s=p,t$ let $n_s$ count the $s$-free
nonunits. At least $k_s=H_{16}(s;2,m)$ have denominators divisible
by $4s^{v_s(M)}$ if $m\ge4$. Put $k_s=0$ otherwise. Use
$n_s\ge\max(2,k_s)$, $n_3\equiv2\pmod3$, $n_5\ge3$, and the
stronger bound \eqref{eq:full-census} at five and seven when $m\ge4$.
The three omission sets are disjoint. When $m<4$ one $p$-free
denominator is divisible by $p^{v_p(M)}t$ and one $t$-free denominator
by $t^{v_t(M)}p$, by principal witnesses.

If $t$ mixes with two or $p$, reserve respectively $4(t-1)$ or
$p(t-1)$ active multiples of $2t$ or $pt$. These avoid the four-row
hits and, when $m>0$, the odd degrees, whose twists are dyadic. For
the odd interaction with $m\le3$ simply reserve the entire $pt$
places of the spherical summand, allowing the vacuum among them.
If $t$ is isolated, first test $t(t-1)$ pure places of floor $t^2$,
allowing them to supply $t$-free places other than four-row hits or
the principal-$p$ witness. Below this count sparse height gives
$M_t=t^{2c-1}$, $c\ge2$, and every nontrivial $t$-part has size
at least $t^2$. Nonminimal activity then reserves $2(t-1)$ pure
$t$-divisible degrees, by primitive balance or a higher orbit.
With fewer than $t(t-1)$ pure labels, minimal activity leaves fewer
than $(t-1)^2$ inactive pure labels. Thus Lemma~\ref{result:packet-idempotent} gives $t$
such degrees at the common denominator $t^c$.
Adding the reciprocal squares of these floors leaves only
$(p,t)=(3,5)$ with minimal activity, $A=8$, and
$m\in\{4,5,7,8,9,10,11,12,13\}$.

Write $M=2^a3^b5^{2c-1}$. The vacuum Weyl congruence gives
$2^a3^b\equiv\pm1\pmod{125}$. Truncate $a,b,c$ at the first
exponents with respective powers exceeding $K=65\,536$, the
last value representing its whole tail. Reserve the vacuum, the
$m$ fours, the five places at $5^c$ and the minimum omission
counts above. For each remaining position enumerate the permitted
divisors $d\le K$ of $M$ with their exact reciprocal squares.
A nonunit degree omits at most one prime. A larger denominator of
floor $F$ contributes in $[0,\max(K+1,F)^{-2}]$, also allowing
padding. The vacuum has positive mass at most $(2^a3^b5^{2c-1})^{-2}$ at the chosen lower endpoints.
If a valuation was truncated, omit its congruence test when necessary
and use the lower endpoint zero for a truncated packet mass.
Enumerate in order within each group, discarding intervals that
cannot contain one. The nine profiles leave no intervals after
$18\,095$ nodes (\texttt{three\_prime\_cutoff.py}).
\end{proof}

\begin{lemma}[Restrictions on mixed twists]
No twist is active at two odd primes. An odd prime mixed with two
has at least $4(p-1)$ active multiples of $2p$, and $p\le5$.
\end{lemma}
\begin{proof}
Suppose two odd primes mix. The spherical summand and the activity
reservations exclude a product $K\ge33$. At $K=21$ the third odd
prime is five. A fourth odd prime gives mass at most
$11/16+1/15+1/7+1/11<1$. Its summand is disjoint as their product
would exceed the rank. A second odd edge, or dyadic interaction at
the third prime, is excluded by the counts and mass (for edges $15,21$
use $12/225+18/441+14/16<1$. For $21$ and dyadic five use
$16/100+18/441+10/16<1$). At $K=21$ reserve its $22$ mixed labels.
Absence of five leaves mass below one even with thirteen fours.
Its presence reserves five labels and four proper witnesses. At
most eleven fours have insufficient mass, and thirteen overshoot.
At twelve all positions are reserved. The four proper five-free
denominators are at least $80$ by Lemma~\ref{result:small-Gauss}.
If any of the $22$ multiples of $21$ exceeds $21$, then
\[
 1\le12/16+1/5+21/21^2+1/42^2+4/80^2+1/420^2<1.
\]
Thus all equal $21$. At equality in Corollary~\ref{result:fiber-squares}
there are exactly two nontrivial ternary fibers meeting seven-activity,
with eight eligible labels apiece. A higher ternary phase would require
three such fibers, a higher seven-phase $21$ labels in one fiber.
Each therefore has six singleton primitive seven-phases and two eligible
seven-neutral labels, all of the same weight. The fiber moments give
$(1-a_1)^2+(1-a_2)^2=8$ with $a_1,a_2\in7\ZZ$, impossible.

At $K=15$ the third odd prime $u$ is isolated. A mixed reservation
and the eighteen multiples of $15$ leave insufficient mass. Thus
$u<44$ by the $u$ pure places. Reserve the vacuum, $m$ fours, these
eighteen places, and either $u$ at $u^2$ or $uj$ at $u$. In the latter
case partition $j(u-j)$ into at least three positive squares $a_i^2$.
The proper witness floors are $uk_i$, with $k_i>1$ least satisfying
$a_ik_i\equiv\pm1\pmod u$. Use Lemma~\ref{result:cross-weyl} at seven.
For $m\ge7$, Lemma~\ref{result:saturation} gives $8\mid M$: reserve three
odd degrees, imposing $8\mid k_i$ on their overlap. The proper
witnesses include an odd degree by gcd one and an even degree since
every odd $u$-free degree contains three and five.
At $j=1$ the prime twist sum vanishes, so odd $m$ requires
$16(1-m/16-1/u+1/(30u))\ge1$. Every remaining upper mass is below one
(\texttt{arithmetic/isolation44.py}).

A dyadic interaction at $p\ge7$ claims at least twenty-four positions
of floor fourteen. A five-packet, its four proper witnesses and two
ternary active labels then give insufficient mass. Otherwise six
positions are active at the other odd primes. For $m\le10$ the bound
is $10/16+24/196+9/36+1/210^2<1$. For $m\ge11$ three odd nonunits
require at least $1+24+6+11+3=45$ places. The count at a mixed prime
is Corollary~\ref{result:mixed}.
\end{proof}

\begin{lemma}
If every active twist has prime-power order and $\xi=\pm1$, the
deficit $1-\xi/M=\sum(1-\langle\theta_i\rangle)/b_i^2$ is a sum of
digits. At odd $p$ a digit $p^{1-2k}$ has $p^k\mid M$ and costs at
least $p-1$ active labels. At a prime denominator the full contribution
is $j/p$. At two each active label contributes $2^{-2k}$ or $2^{1-2k}$.
\end{lemma}
\begin{proof}
Local splitting makes an active denominator a power of its twist prime.
The primitive $p$-phases have equal weight $W$ and deficit $pW$. Use
the least numerous fiber to charge $p-1$ labels per term. Higher
phases have mean zero and occur in equal-degree packets divisible
by $p$. The dyadic formula follows from the root trace.
\end{proof}
In the resulting exact recursion, if the current nonzero remainder
is $A$, at most $s$ digits remain, and their largest size is $1/u$,
then $|A|\le s/u+u^{-1/2}$, since the final error is $\pm1/M$.
Each choice imposes its displayed divisor on $M$. A zero intermediate remainder would leave a positive sum of digits
equal to $1/M$. At an odd $p\mid M$ the sum of the remaining $p$-digits
would then have valuation $-v_p(M)$, since every other digit is
$p$-integral. Its reduced denominator would be $M_p$, giving a sum
at least $1/M_p>1/M$.

\begin{lemma}
Exactly four primes divide $M$.
\end{lemma}
\begin{proof}
Each odd $p$ needs $p$ pure labels if isolated and $4(p-1)$ if mixed.
Five odd primes require at least $39$ places and only the prime three could
mix with two. The upper mass with all isolated is
$1/27+1/5+1/7+1/11+1/13+4/16+1/2310^2<1$.
Mixing three would require $8+5+7+11+13+1=45$ labels. Larger choices
only reduce the bound.

For four odd primes reserve at each $p$ one of $(p,p^2)$,
$(4(p-1),2p)$, or $(p,p)$ according as $p$ is absent and isolated,
mixed, or present. A present denominator contributes at least three
proper witnesses at $2p$, or four at $20$ if $p=5$. For $m\ge9$
reserve three odd labels of floor eight, allowing overlap with the
witnesses by taking coordinatewise maxima of their decreasing floors.
Use $M\ge2310$ and default floor six. Enumeration gives $49$
profiles, of which the $38$ all-primary ones are excluded by the
preceding deficit recursion. Here at least one prime denominator
is present. Subtract its $j/p$ and $(m-\langle T_4\rangle)/16$,
and charge the vacuum, fours, prime classes and largest proper
witness quota before enumerating the remaining digits.

In the other eleven profiles eight reserved denominators are multiples
of six supported on two and three. Reserve also $p$ pure powers
for each absent isolated prime and $p$ copies of each present prime.
At most eight positions remain. A present prime has one packet.
This holds at five, and another packet at $p\ge7$ would leave fewer
than the three required proper degrees. On the remaining positions
its Weyl coefficients satisfy
\[
 \sum_i a_{pi}^2=p-1,\qquad \sum_i a_{pi}/b_i=\epsilon/M,\qquad
 \sum_i a_{pi}a_{ti}=-1\quad(p\ne t).
\]
Indeed $(W_p-I)e_0$ and $(W_t-I)e_0$ are orthogonal since no summand
is nontrivial at both primes. Deleting the vacuum gives the last
identity. If $p\nmid b_i$ then $a_{pi}=0$. Otherwise it is the least
absolute inverse of $b_i/p$ modulo $p$.

Enumerate the coefficient magnitudes with the indicated norms, at
least three nonzero entries per row and odd pairwise dot products.
Within each resulting column type try the compatible denominators
through $256$ in increasing order, allowing repetitions. A larger
denominator contributes mass in $[0,257^{-2}]$ and at most
$|a_{pi}|/257$ in modulus to the signed sum. The zero endpoints
allow padding to rank $44$. For the vacuum use
$M\ge4(\prod_{p\mid M,\,p>2}p)\prod_{t}t$, where $t$ runs over
the absent isolated primes. The mass test leaves eighteen intervals. The signed sums exclude all eighteen (\texttt{support\_five.py}).
\end{proof}

\begin{proposition}[Multiplicity at denominator four]\label{result:thirteen-34}
With $m\ge9$ a twist is active at two and three, and every other
odd prime is isolated. At $m=9,10,11,12,13$ the numbers of admissible
mixed profiles are $12,11,9,6,1$. Through rank $41$, $m\le12$.
Through $40$ the odd primes are $3,5,7$, with five and seven
isolated. For $m=9,10,11$ in the latter range, absence of denominator
five forces seven. At eleven both cannot occur, and at twelve seven
is absent.
\end{proposition}
\begin{proof}
Here $M/4$ is even and $m\le13$. The vacuum, $m$ fours and at least
three odd hits avoid all odd activity. At an isolated absent $p$ reserve
$(p,p^2)$, at a mixed $p$ reserve $(4(p-1),2p)$, and at a present
$p$ reserve its prime packet. For every odd $p$ record a divisor $P_p$
of the denominator of a $p$-free degree: use $p$ for a mixed or present
$p$, $9$ for absent isolated three, $p^3$ for other absent isolated
primes. In the last case the sparse-height bound applies. Indeed the
vacuum, fours, odd labels and other prime activity leave fewer than
$p(p-1)$ positions at $p\ge7$. At absent isolated five, twenty active
labels give mass at most $20/625+10/16+3/64+7/49+3/36+1/210^2<1$.
At eighteen, three further five-free positions have floor $125$
since sparse height applies below twenty. Eleven fours still give
insufficient mass. Twelve require precisely three remaining positions,
three pure ternary and seven pure seven labels on $2,3,5,7$. Those
three remaining would be both five-free and odd. Without seven the
mass is too small, while its packet needs an odd seven-free witness,
necessarily then of prime-power degree. Thirteen fours cannot fit.
Hence at most sixteen are five-active. An odd prime $\ge23$ exceeds
the rank or has insufficient mass already at $23$.

Group the remaining labels by the odd primes $J$ omitted by their
degree. A denominator floor is $c\prod_{p\in J}P_p$, with $c=8,4,1$
for an odd degree, even four-row hit, and other even degree. The
possible $|J|$ are at most one for odd degrees and two otherwise.
Impose the weighted and unweighted bounds of Lemma~\ref{result:joint-projection}
on the hits (in the latter allow either common parity of their row
sums). At a present prime the proper free degrees omit each other
prime. Of the four five-witnesses exactly one is odd, by parity and
the two distinct even witnesses omitting the other odd primes.
At least $\lceil m/(16-m)\rceil$ are hits by Lemma~\ref{result:cross-weyl}.
At twelve use also Lemmas~\ref{result:cross-weyl} and~\ref{result:small-Gauss}.
Thirteen fours exclude five by mass and seven by the impossibility
of $T_4=0$ with odd cardinality (one packet leaves $5/112$, two
overshoot). Set $M\ge8\prod P_p$ in the vacuum slot.
These are the finite omission tests. Retain every nonnegative
profile for which the mass can reach one.

The omission sums give $41$ mixed signatures. Only two mix five, both
on $2,3,5,7$ with three isolated and seven present. Their omission
lists at ten and eleven fours number five and one. Each contains three
odd witnesses omitting $3,5,7$, respectively, and two even hits
omitting $\{3,7\},\{5,7\}$. At eleven the sole other witness omits
$\{3,5\}$. All positions are accounted for, with seven sevens but only
three proper seven-free labels, contrary to Lemma~\ref{result:cross-weyl}.
At ten, after the vacuum, fours, sixteen mixed labels, seven sevens and
the five witnesses, at most five positions can be ternary-active.
Lemma~\ref{result:activity-height-gap} gives $27\mid M$ and excludes
six. Reserving three pure ternary labels gives the contradiction
\[
 1\le\frac{10}{16}+\frac{16}{100}+\frac3{81}+\frac17+\frac2{64}
 +\sum_{b\in\{216,40,56,756,140,7560\}} b^{-2}<1.
\]

For the all-primary case reserve at least three neutral labels by the
four singleton twist components (Corollary~\ref{result:mixed-neutral-rank}).
Take the maximum with the proper-witness quota. Five and another prime
packet require at least six proper labels, since two at the latter
must omit two and the remaining odd prime separately and cannot be
five-free. Subtract prime contributions $j/p$ and $(m-t)/16$,
$t=\langle T_4\rangle$. The digit recursion with these reservations
leaves, from $46$ profiles, just $M=136080,15482880$, both with nine
fours and one packet at five and seven (a second seven-packet
overshoots). Their proper signed unions number $13,39$. Both prime
twist sums vanish and the residual mass gives $|t|\le1$. Enumerate
the ternary digits $3/b^2$ at $b=9,27,\ldots\mid M$, each costing
two places, and all dyadic splittings. This leaves $0,218$ profiles.
For ternary deficit $G$, $h$ open positions and $v$ nines,
$2G/3\le v/81+(h-v)/729$. Beyond the nines, all remaining degrees
are multiples of $35$ dividing $M$. Divide the residual square sum
by $35^2$ and test all divisor squares with zero padding. None of the
$5092$ targets on at most five slots occurs. The omission sums,
deficit recursion and terminal square sums are the finite comparisons
implemented by \texttt{preflight.py} in the census source.
Using the actual rank in the same sums gives the asserted smaller-rank
restrictions (ten mixed profiles through rank forty).
\end{proof}

\section{Odd dimension}\label{sec:odd-new}
Write $D=qM^2$ with $q$ squarefree. In odd dimension every degree and every
denominator $b_X=M/d_X$ is odd. The nonunits occur in dual pairs and
$r\equiv q\pmod8$ (Lemma~\ref{alg:odd-parity}). Let $m_b$ be the multiplicity of $b$.
We use dual pairs as the units of counting unless a count is called a count of labels.

\begin{lemma}[The diagonal entry and the twist]\label{odd:diagonal-twist}
The field of $X$ contains its twist. More precisely,
$\chi_X(X^*)\equiv\theta_X^2\pmod2$.
\end{lemma}
\begin{proof}
In the $(X,X^*)$ entry of $STS=p^+T^{-1}ST^{-1}$, the terms indexed by a
dual pair agree. Modulo two the vacuum is the only term and $p^+\equiv1$ by
the same pairing. The stated congruence follows since $d_X$ is odd.
An automorphism fixing its left side fixes $\theta_X$: odd-order roots
of unity inject under reduction at two.
\end{proof}

At fixed complementary phases the mass identities
(Lemma~\ref{block:moments}) hold in pair units for every fiber excluding the
vacuum. For a set of primes $p$ of even valuation choose disjoint sets of
$p-1$ active phase tuples, by the packet bound for each subset and Hall's
theorem. At a prime $t\mid q$ each such tuple costs at least
$C(t)=t$ pairs if $t\not\equiv1\pmod {16}$ and $C(t)=(t+3)/2$ otherwise.
Indeed, a twist of order $t^2$ already gives $t(t-1)/2$ phases. At level $t$,
put $h=(t-1)/2$. The primitive count is $hk$, and the mass identity,
reduced modulo eight using odd squared degrees, gives $hk\equiv2h n_0\pmod8$,
where $n_0>0$ is the neutral count. Hence $k$ is even if $8\nmid h$.
Otherwise use Lemma~\ref{block:moments}. The same bounds hold in labels if
the fiber contains the vacuum.

For any nonunit $X$ write $z_Y=2\Re\chi_X(Y)$ on the nonunit pairs.
Orthogonality of $X,X^*$, orthogonality to the vacuum, and the sum of their
vacuum modular relations give, respectively,
\begin{equation}\label{odd:pair}
 \sum z_Y^2=f_X-2,\qquad \sum d_Y z_Y=-1,\qquad
 \sqrt D\le1+\sum d_Y|z_Y|.
\end{equation}
A whole Galois orbit of nonzero algebraic integers spends at least its
cardinality in the first sum, by the arithmetic--geometric mean inequality.
Its sum of absolute values is at most its squared cost, by Cauchy--Schwarz.

\subsection{Nonsquare dimension}
For the trace estimates put $D=3M^2$, $n=M/3$. A denominator-three row has
nonreal field in $\QQ(\zeta_9)$: degree $18$ in $\QQ(\zeta_{27})$ is
impossible by Lemma~\ref{small} (at most two nonzero-mean
orbits). Each hit has degree dividing $n$, by principal-block height.

\begin{lemma}\label{odd:27}
On three-prime support, denominator three is impossible in $3M^2$.
\end{lemma}
\begin{proof}
Write $3,p,l$ for the primes. Every $s^2\mid M$. The principal $s$-block
has a nonunit degree divisible by $s$, since its degrees have gcd one
and any nonunit degree has at least two primes. Its denominators are
$s$-multiples. In \eqref{odd:pair}, each nonzero orbit at a three-divisible
degree costs at least nine. Its trace is either an integer divisible by
three or a real cubic cyclotomic integer reducing to zero at three. In
the second case the trace of its square is a multiple of three at least
$3\cdot3^{2/3}$. A $p$-free and an $l$-free hit with nonzero trace are
required, in distinct such orbits, leaving at most seven of the norm $25$.
The remaining three-free hits have degrees at most $n/3$ and the two
required orbits have degrees at most $n/p^2,n/l^2$. Thus
$3\sqrt3\,n\le1+7n/3+25n/\min(p,l)^2$, which is impossible.
\end{proof}

At a ramified $t$, both degrees of a defect-one block are odd, so in
Lemma~\ref{block:defectone} their squared ratio is an odd square $a\mid t-1$.
The alternative $2a$ has positive valuation at two. If $t-1$ contains no
nontrivial odd square factor, the block consists of $t$ equal degrees, one
$t$-neutral. Removing all such ramified factors at a fixed complementary
phase divides each multiplicity prime to the ramified denominators by their
product and leaves fields in $\QQ(\zeta_{b^2})$. In particular at $b=p^a$,
the residual multiplicity is divisible by $2^{v_2(p-1)}$, since the
stabilizer in the cyclic group has odd order. Denominator one is excluded.

At an unramified prime $p$, an order $p^j$ square class in a complete class
of denominator $b$ costs a multiple of $p^{j-1}v$, where $v=p-1$ for
$t\mid b$ and $v=t(p-1)$ otherwise ($q=t$ prime). This follows by square
covariance, duality, and the equal $t$-blocks. At $b=p$ and $r<10v$ the two
primitive classes have equal multiplicity. Their discrepancy would cost at
least $v/p^2$ beyond $p$, at floors $3p$, hence $9v$ further labels. The
possible activity counts at $b=p$ are consequently $2vh+pvk$, $h,k\ge0$.

\paragraph{Effective rows at an unramified prime.}

Suppose $q=t$ is prime, $t-1$ has no nontrivial odd square divisor and
$p\ne t$ occurs as denominator. Lemma~\ref{result:height} gives $p\Vert M$.
Suppose all $p$-activity is at $p$. Write $m=m_p/t$, $n=M/p$. A common
even number $j$ of blocks occurs at each primitive phase, by the mass
identity and duality, and $0<j<p$ by the multiplicity bound. Then
\begin{equation}\label{eq:effective-weyl-bound}
 m\ne p(p-1),\qquad p|pj-m|\le j(p-j)+(p+1)j/n.
\end{equation}
Indeed apply Lemma~\ref{lem:congruent-rows} to the effective rows for the
first assertion. For the second, sum their active rows to $L$ and let
$w=W_pe_0$. The vacuum Bruhat identity at exponent $k\in\mathbb F_p^\times$ reads
$\theta_{Y,p}^{-1/k}w_Y-\delta_{Y0}
=p^{-2}\sum_{B\ {\rm active}}(\theta_{B,p}^k-1)v_B(Y)$.
Averaging gives $p(w_Y-\delta_{Y0})=L_Y$ on active labels and
$-L_Y/(p-1)$ elsewhere. Thus $A=pw$ is supported on the $t$-neutral
labels, $A_0=p-j$, $A_Y\equiv-jd_Y\pmod p$ there.
Indeed $w$ is rational, integral away from $p$ by the displayed identity,
and $b_Yw_Y\in\ZZ$, $p^2\nmid b_Y$. Each neutral singleton $X$ in the
$m$ blocks has $|w_X|=1/p$, since $W_t^{-1}$ transfers its vacuum
projection isometrically to this line. The congruence fixes its common
sign $\delta$ in $A$. For the proper support put $a_Y=\delta A_Y$,
$k_Y=n/d_Y>1$. Then $k_Ya_Y\equiv1\pmod p$.
From $\sum d_Y L_Y=0$ obtain $\sum d_Y A_Y=p+pnj\delta$.
Subtracting the vacuum and the singletons here and in $\sum A_Y^2=p^2$
gives \eqref{eq:weyl-cost} with $z=pj-m$, hence the bound.

\begin{lemma}\label{odd.Two-neutral}
If $p\ge5$, $j=2$, and $2n>p+1$ in the effective case, then $m=2p$, $p\equiv3\pmod4$, $p\ge11$.
The proper coefficients, one per dual pair, have squared sum $p-2$
with exactly one odd term.
\end{lemma}
\begin{proof}
For the ordinary rows $m=2p$ by Proposition~\ref{result.One-packet}.
For effective rows, $m$ is even by duality and
$p|2p-m|\le2(p-2)+2(p+1)/n<2p$ by \eqref{eq:effective-weyl-bound}.
The effective rows on the $t$-neutral charge-even lines are scalar
multiples of $W_pv_X$, equal on dual blocks (Proposition~\ref{result:effective}).
By Lemma~\ref{result:steinberg-compression} the neutral compression is $-gg^*/(2p)$, where $g_0=-2e$,
$g_Y=\sqrt2a_Y$ and $e$ is the common sign. Consequently the neutral
row has pair traces $-ea_Y$ on proper pairs, $-e$ on active pairs
(by \eqref{eq:packet-trace}), and $e(p-1)$ on its own pair.
Its nonreal quadratic field in $\QQ(\zeta_{p^2})$ is $\QQ(\sqrt{-p})$,
so $p\equiv3\pmod4$. Write the entries as $(z_Y+c_Y\sqrt{-p})/2$.
Row orthogonality gives $\sum c_Y^2=p$. Since $p-1$ active pairs
already have odd trace and $\sum a_Y^2=p-2$ is odd, exactly one
proper coefficient is odd. At least three distinct proper degrees
are needed (gcd one, pairwise noncoprime), hence $p-2\ge1+4+4$.
\end{proof}

\paragraph{The ternary class.}
Assume $D=3M^2$ and let $F$ be the three-free labels including the unit.
Every denominator-three row meets $F$, so $|F|\le27$. Deletion and
the dimension equation make $|F|$ an odd multiple of three at least nine.
At equality $|F|=27$ the rows vanish elsewhere. If $9\mid M$, Bessel
applied to $(d_Y\theta_Y)_{Y\in F}$ and the modular relation gives
$m_3/9\le\sum_F b_Y^{-2}\le27/81$, so $|F|\le21$ for $m_3>3$.

\begin{lemma}\label{result:ternary-cap}
On four-prime support with $9\mid M$, a nonquadratic denominator-three
row has order-nine twist and $m_3\le12$. If at least one of $27\mid M$, $m_9\le2$, $m_3>3$ holds, a
nonempty three-class forces a nonunit pair of degree coprime to $3p$ for
every other prime $p$. In these ranges the quadratic three-rows have
equally many occurrences of each cube-root phase.
\end{lemma}
\begin{proof}
For a nonquadratic row $X$, put
$u=\chi_X(X^*)/\theta_X^2$. It is one at two and zero at three.
Both $v=(u-1)/2$ and $v+1$ are nonzero integers, so
$C(u)=2C(v)+2C(v+1)-1\ge3$. The six conjugate positions of $X^*$ and
$F$ are disjoint, costing at least $18+9=27$. Equality forces $|F|=9$,
$u=\pm\sqrt{-3}$ (both $v,v+1$ are roots of unity) and all entries on
$F$ roots of unity. Its mean over the cubic subgroup has squared norm nine
by row orthogonality. If $\theta_X^3=1$ this mean retains the entire
orbit cost $18$. Hence $\theta_X$ has order nine, and averaging kills
the orbit entries. The entries on $F$ are
therefore sixth roots, already fixed by the subgroup. The mean is
orthogonal to its dual: $\sum_F\chi_X(Y)^2=0$. Thus exactly one free
pair has doubled trace $2e_i$, the others single traces $-e_i$,
where $e_i=\pm1$ on each pair is fixed by reduction at three.
They satisfy $1+\sum_F' d_i z_i=0$, the prime indicating one term per
nonunit pair.

For a quadratic row let $f=(|F|-1)/2\in\{4,7,10,13\}$ and
$h=\min(v_3(M),3)$. The free trace magnitudes $v_i$ are positive
integers prime to three. At three-divisible degrees there are $j\le2$
nonzero traces, all $\pm3$, so $\sum v_i^2=25-9j$. Let $\ell$ of these
$j$ occur at degree $n$, with signed sum $A$. At $h=2$ exactly
$k=m_9/2$ free pairs have degree $n/3$. Call their indices $J$.
At $h=3$ put $J=\varnothing$, $k=0$. The remaining free degrees are
at most $n/15,n/9$, respectively, and remaining three-divisible hits
at most $n/5,n/3$. Write $B=\sum_{i\notin J}v_i$ and
\[
 H=\begin{cases}B/15+3(j-\ell)/5&h=2,\\
 B/9+j-\ell&h=3.\end{cases}
\]
The signed traces on $J$ agree modulo three. Write their sum as $L$,
so its two possible values follow from the magnitudes.
By \eqref{odd:pair},
\[
 |A+L/3|\le H+1/15,\qquad
 3\sqrt3\le3\ell+\tfrac13\sum_{i\in J}v_i+H+1/15.
\]
The proper hits must omit all three unramified primes. Hence
$f-k+2(j-\ell)\ge3$. If $A+L/3=0$, at least two hits outside these degrees remain
to contribute $-1$. Checking the partitions of $25-9j$ into $f$ squares
and their subsets $J$ gives $(h,k,\ell,A)=(2,0,2,0),(2,1,2,0),(3,0,2,0)$.
Here $f=13$ is excluded when $k\ge2$, since in that range $m_3>3$ by
hypothesis. Thus $9\Vert M$, $m_9\ge4$, $m_3>3$ permits no quadratic row.

The two opposite degree-$n$ hits are each Galois-stable pairs,
hence quadratic. The remaining weighted absolute trace mass is at
most $13n/15$ from the residual norm seven.
All quadratic rows have cube-root twists. Multiplier four fixes their
field and raises their twist to the sixteenth power. If $\theta_+,\theta_-$
are the twists of the two hits, their difference is within $1/3$ of
$\pm\sqrt{-3}\theta_X^{-1}$ by the modular relation
($p^+=\pm M\sqrt{-3}$, by the odd conductor and the two-power Gauss
phase). Differences of cube roots force equality. A pair of phase
exponent $a$ thus sees one pair each of exponents $1-a,2-a$ modulo
three. Symmetry gives equal numbers at the three phases.

In all cases considered each row gives $1+\sum_F' d_i z_i=0$, so a
three-free pair omitting each other prime exists. If $|F|>9$, all rows
are quadratic and squared free norm seven forces seven pairs with
constant signs. For $|F|=9$, every row has one double and three singles.
The doubled index satisfies $3e_i d_i=\sum_j e_jd_j-1$. Since three
distinct proper degrees are required, only two indices can be doubled.
Two normalized symmetric pair columns consequently have constant entries
of magnitude $1/\sqrt2$ on these rows. Together with the vacuum column,
Bessel gives $2m_3\le27$. Duality yields $m_3\le12$. This reasoning
also covers the cap outside the indicated ranges, where $m_3\le3$.
\end{proof}

\paragraph{Omissions and edges.}
If a five-power codegree $f$ has $m>f/2$ copies, every five-free nonunit
field $K_Y$ contains $\QQ(\zeta_5)$. Indeed its subfield $E$
generated by the class entries at $Y$ lies in $\QQ(\zeta_f)$ and
$[K_Y:\QQ](m/f+1/f_Y)\le[E:\QQ]$ by
Lemma~\ref{result:field-rank}. Thus $K_Y=E$, and the assertion follows
by nonreality and cyclicity. Two such
orbits are needed by deletion. Eight labels would give
$1+4a^2+4b^2\equiv0\pmod5$ with $5\nmid ab$, impossible. Thus at
least twelve occur. More generally, suppose $p\equiv1\pmod4$ and
$m$ rows of $p$-power codegree $f$ satisfy $mp^2>f$. If all their
$q$-free nonunit hits have degree and denominator divisible by $p$, there
are at least four hits. Otherwise their single dual pair has imaginary
quadratic field disjoint from $\QQ(\zeta_f)$, and the $m$ entries are
nonzero integers divisible by $p$, exceeding the column norm.

In $3M^2$ joint activity at $p,t\ge5$ costs at least
$12(p-1)(t-1)$ jointly active labels and $2(p-1),2(t-1)$ singly active
labels. In primitive phases the counts per phase in pairs agree modulo
eight by equal masses, and are even by the two-moment parity. At fixed
joint phase a higher three-phase then costs six pairs. At level three
the positive even neutral and active counts satisfy
$n_1\equiv2n_0\pmod8$, so again cost six. There are $(p-1)(t-1)$
joint phases. A higher $p$-order supplies $p(p-1)/2$ phases, each with
at least $t-1$ active $t$-phases costing six labels each. The second
moment requires an inactive pair over every occupied complementary
phase, giving the singly active counts.

\paragraph{Prime support.}
If $q>1$ and $r<217$,
square-free exclusion gives $9q<r$ and $q\le23$. For $q=3s=15,21$ an
unramified $c$ gives $r\ge9q+4s(c-1)(1-9/c^2)>216$
($c\ge7,5$, respectively). At $q=19,23$ use
$r\ge9q+2C(q)(p-1)(1-9/p^2)$ for an unramified $p\ge5$.
Thus $q=t\le17$ is prime. All $t$-blocks at an unramified prime
denominator have equal degrees.

On three-prime support each prime square divides $M$ by the principal
block gcd argument of Lemma~\ref{odd:27}, excluding all unramified
prime denominators by the height theorem. For $t=3$ Lemma~\ref{odd:27}
now gives $r\ge81t$, and for $t\ge7$ one gets $r\ge49t$.
For $t=5$ we have $r\le213$, so $m_5>62$ by the mass bound
$62/25+151/81<5$. Reserve twelve five-free labels at floor $25$.
At each other prime $p$ a $p$-free hit of a five-row has five-divisible
degree and denominator divisible by $5p^2$. The preceding omission bound
reserves four such labels. These sets and the active sectors are disjoint.
Reserving twenty active labels at nine and sixty at twenty-one gives
$112/25+20/81+60/441+12/625+4/45^2+4/245^2+11025^{-2}<5$.

With at least four primes, whenever $j=2$ one has $2t(p-1)\le r$ and $n\ge3\cdot5\cdot7$, so Lemma~\ref{odd.Two-neutral} applies. The inequality $C(t)(1+2\sum_{p\ne t}(p-1))\le r$
leaves $t=3,5$. For five the possible supports are $3,5,7,11$ and
$3,5,7,13$, with disjoint unramified activities (a $3,7$ edge costs
$120$, apart from one hundred at the last prime). The effective
multiplicity at denominator three is even, at most seven and unequal to
six by Lemma~\ref{lem:congruent-rows}, giving $m_3\le20$.
Concentration at three forces $m=4,j=2$ in \eqref{eq:effective-weyl-bound},
hence $6\le2+8/n$. Concentration at seven excludes $j=2$ by
Lemma~\ref{odd.Two-neutral}. $j\ge4$ costs $120$ labels beside
twenty and one hundred. At three an imbalance costs one hundred
further labels besides the other $160$. At seven use the first-active
denominator bound. Removing the balanced equal quinary blocks therefore
requires the full $20,60$ residual activity at floors nine and
twenty-one. The mass is at most
$20/9+12/25+20/81+60/441+100/121+1155^{-2}<5$.

Hence $D=3M^2$, $r\le211$. Four unramified primes require precisely
$P=\{5,7,11,13\}$. The edge bounds above make their activities disjoint,
costing at least $A_p=6(p-1)$ apiece. At $p=5,7,13$, $j=2$ is impossible
by Lemma~\ref{odd.Two-neutral}, while $j\ge4$ plus the other costs and the
vacuum gives $217$ labels. An imbalance at $p$ already costs
$30(p-1)+\sum_{s\ne p}A_s>211$. Removing balanced ternary blocks and
using the ramified mass identity therefore reserves $A_p$ labels at $3p$
at these three primes. The upper mass is
$18/9+24/15^2+36/21^2+60/11^2+72/39^2+15015^{-2}<3$.
There remain three unramified primes, with $3+\sum A_p\le211$ (fourteen
supports). Their activities are again disjoint and their primitive phase
classes balanced at prime denominators. If activity persists after
deleting the ternary blocks at $p$, it costs $A_p$ labels at floor $3p$.
Otherwise use \eqref{eq:effective-weyl-bound} and Lemma~\ref{odd.Two-neutral}.

\paragraph{Multiplicity and mass.}
At most one unramified prime $p$ is a denominator, with $p\Vert M$ and
the other two heights at least two (Lemma~\ref{block:defectone}).
Bessel on its effective rows, subtracting the vacuum row from its modular
image, gives
\[
 m_p/(3p^2)\le\sum_{\theta_{Y,3}=1,\ \theta_{Y,p}\ne1}
                     |\theta_{Y,p}-1|^2/b_Y^2.
\]
The averaged squared factor is at most $2p/(p-1)$ and at most a third
of the mass at a nontrivial $p$-phase is three-neutral. With no activity
at $p$ itself this requires $9(p-1)m_p/(2p)$ active labels at $3p$ or
above. For general $t$ replace both factors of three by $t$.

Two joint restrictions on $m_3,m_5$ are needed. Choose a three-free
nonunit hit $Y$ in an effective five-row. It is also five-free and its
orbit meets every row in both families. At $(m_3,m_5)\ge(8,36)$ their
masses sum to more than $3/4$. Their fields are disjoint, so
Lemma~\ref{result:field-rank} forces one family rational on $Y$ and
$|G Y|\le\max(27/8,25/12)<4$. On this pair the five-values are nonzero
integers, and every ternary pair trace is nonzero modulo three. The
symmetric pair column costs $24/25+m_3/54>1$. At $(m_3,m_5)\ge(6,48)$
the same test forces the three-family rational since the five-family
alone exceeds half. Now $|G Y|\le27/6$ and the five-family forces
$\QQ(\zeta_5)\subset K_Y$ by the pair column bound, so equality of the
fields holds. Every five-trace in this quartic field is $4d_Y\ne0$
modulo five. Projection on the constant orbit line costs at least
$4m_3/27+m_5/300>1$.

The finite sum checks $m_3,m_9$ and the three prime multiplicities
($m_5\le48$ by the effective bound, divisibility by twelve and the
exclusion of effective multiplicity twenty). A remaining nonvacuum
denominator has floor $15$. After the preceding inequalities every
survivor has a pure three-power row. All three-free labels are then
inactive at the other primes. The deletion bound requires eight
three-free nonunits. A prime of height one forces a dual pair omitting
it and each other prime by principal-block height and gcd one.
Different omitted pairs are disjoint. The free trace equation of
Lemma~\ref{result:ternary-cap} also forces the omissions $\{3,p\}$ in its
range. A pair omitting $p,s\in P$ is inactive at every unramified
prime by Lemma~\ref{result:spherical}. If its denominator is coprime to
three it belongs to six equal-degree labels, so fewer than four spare slots force a two-label witness to have denominator divisible by three. Use the full known powers
of the omitted primes in these floors and the vacuum bound. At
$9\Vert M$ the unrecorded three-free labels have floor $45$. If also
$p\Vert M$, a residual $p$-active label has denominator exactly $3p$
(its support lies in $\{3,p\}$ and at $9p$ it would be three-free).
If $3\Vert M$, a residual active denominator at $p$ is a pure
$p$-power $\ge p^2$, since $3\mid b$ makes its degree three-free.
This excludes residual activity for $p\Vert M$. In this case count
the recorded three-free labels and their reservations against the
bound $|F|\le27$. Once exhausted, all further denominators are
prime to three and at least $25$. Credit each position once.

Enumerating the multiplicities at $3,9$ and the three unramified primes
with total at most $210$ and mass below three leaves four valuation branches.
Each has $P=\{5,7,11\}$, $M_0=27\cdot5\cdot49\cdot121\mid M$ and
$(m_3,m_5,m_7,m_9,m_{11})=(12,24,0,k,0)$, $k=36,38,40,42$.
Six disjoint unrecorded pairs outside the activity have denominators
divisible respectively by $135,1323,3267,27,245,605$, from the three
omissions $\{3,p\}$, a fourth three-free pair and the principal five-block.
The three activity reservations cost $6(p-1)$ each and leave at most
$42-k\le6$ spare labels. Higher $p$-order costs at least $3p(p-1)$, so
all phases are primitive. Their counts after deleting the five-blocks
are equal. The two square classes agree modulo eight in each phase
and an imbalance needs $4(p-1)$ extra labels. A common increase needs
$2(p-1)$ by duality. Hence every primitive phase has six labels.

These six labels have equal degrees. A three-phase of order at least
nine exhausts the three pairs by square covariance. Otherwise the
ramified mass identity modulo eight gives two primitive pairs and
one neutral, with $d_1^2+d_2^2=2d_0^2$. Reducing at three after division
by the lowest common three-power shows that all three valuations agree.
Their denominators are supported on $\{3,p\}$ by
Lemma~\ref{result:spherical}. Thus their degrees differ by powers of
$p$. Division by the smallest and reduction at $p$ shows they coincide.
The complement of the removed five-blocks is invariant, so the mass
identity applies there as well. Equality of primitive phase masses
gives a common denominator $B_p=3^i p^j\ge3p$ at each $p$.

Thus the activity mass is $\sum 6(p-1)/B_p^2$. Check $B_p\le500$ and
bound each larger reciprocal square in $[0,501^{-2}]$, reserving the
six pairs, $M_0^{-2}$ for the unit and $1/225$ per spare position.
Only $k=36$, $(B_5,B_7,B_{11})=(15,21,33)$ remains, forcing $r=211$.
There the residual identity is
$5023/266805=M^{-2}+\sum_{i=1}^{9}2/b_i^2$. For residual $E$ with
$h$ pairs left the next sorted denominator satisfies
$\sqrt{2/E}<b\le\sqrt{(2h+1)/E}$. Try all supported $b$ in this interval
with $v_5(b)\le1$, all possible assignments of the six divisors, and at
a leaf require that $M$ be a multiple of $M_0$ and all $b$.
The $72$ nodes have four leaves, none admissible.

\subsection{Square dimension}
Now $q=1$, $r\equiv1\pmod8$. We shall bound the number of labels in
each twist component and compare the remaining ordinary mass with the signed Gauss sum.

\begin{lemma}\label{odd:nine-support}
For $27\mid M$ or $m_9\ge9$, a nine-row is supported on $9\mid b_Y$.
\end{lemma}
\begin{proof}
At a hit $3\mid b_Y$. If $b_Y=3k$, $3\nmid k$, put
$u=\chi_X(Y)/3=\chi_Y(X)/k$, integral by symmetry and lcm integrality.
For $9\Vert M$ every nine-row meets $Y$ by its three-free degree
and principality. The column norm, averaged over Galois, gives $9m_9<81$.
For $27\mid M$, $u$ vanishes above three and belongs to
$\QQ(\zeta_{81})\cap\QQ(\zeta_{9k^2})=\QQ(\zeta_9)$.
Set $h=[\QQ(u):\QQ]$. Its squared trace
$T=\Tr(u\bar u)$ is a positive multiple of three, at least $h3^{2/h}$
by its norm and total ramification, and strictly below nine by summing
the column over the conjugate rows. Thus $h=2$, $T=6$, and $X$ has
orbit of size two (another copy would exceed the budget).
Write its entries $(z+c\sqrt{-3})/2$. Orthogonality of the two rows
gives $\sum z^2=79$ and $\sum c^2=27$ on nonunit pairs. For
$\chi_X(Y)=3u$ of norm $27$, the Eisenstein integers of norm three give
$|z|=9$ or $|c|=6$, impossible.
\end{proof}

\begin{lemma}[The prime support of a character field]\label{odd:field-component}
Each nonunit field is supported on one twist component. At an isolated
prime denominator $p$ let $d=2^{v_2(p-1)}$ and let $m,j$ be as
in Lemma~\ref{result:weyl}. Then $d\mid j\mid p-1$, $j<p-1$, and $(p-1)/d\ge5$.
\end{lemma}
\begin{proof}
For a component $J$, let $E$ be the $J$-local span of its nontrivial
phase spaces. The complementary upper and lower unipotents are scalar
on them by the local tensor decomposition, so the complementary tori
are scalar and, as unsigned permutations, fix them pointwise. They
thus fix $E$. The $J$-group fixes $E^\perp$. For inertia permutations
$P,Q$ from opposite sides, $(P-I)(Q-I)=0$. Two permutations moving
$Y$ would give a positive $Y$-coordinate on $e_Y$, proving localization.
At a prime denominator the level is $p$. A moved row field is nonreal
in $\QQ(\zeta_p)$, hence its orbit size is divisible by $d$.
Proposition~\ref{result.One-packet} makes the neutral class one orbit of size
$j$, a proper divisor of $p-1$. The torus deficit is at most $2j-2$ by Lemma~\ref{result:neutral-inertia}.

If $(p-1)/d=3$, necessarily $j=d$ and the neutral orbit consumes $d-1$ of this deficit. The remaining deficit $d-1$
permits only one further moved orbit, of size $d$. The coefficients
at fixed labels are divisible by $d$ by Proposition~\ref{result.One-packet} at orbit length one.
There are at least three distinct proper degrees. Their squared norm
is at least $6d^2$ with only fixed pairs, and $d+4d^2$ otherwise,
against $d(p-d)$. This exceeds the budget for $(p-1)/d=3$. The case $(p-1)/d=1$ is already excluded by $j<p-1$.
\end{proof}

\paragraph{Parity of local Weyl columns.}
Let $C_J=W_J^2$. Then $(W_J)_{Y0}^2\equiv(C_J)_{Y,Y^*}\pmod2$.
Indeed $W_J$ is symmetric. Its inverse is an upper unipotent, its
lower Weyl conjugate, and the same upper unipotent. The middle factor
$sUs^{-1}$ is symmetric since $U$ commutes with $s^2$.
All are integral above two and commute with duality, so dual terms in
the square cancel modulo two. For a component $J$ the local sign
therefore detects precisely the fields on $J$. In the component norm
identity \eqref{eq:component-norm}, the squared sum on nonunit pairs
equals $g_J-g_J^2/2=2t(1-t)$ with $g_J=2t$, $t\in\ZZ_{(2)}$. It
vanishes modulo four. Thus the component's population is divisible by eight.

\begin{lemma}
At an isolated $p$, let $A$ count active labels and $Z$ moved neutral
labels. Put $d=2^{v_2(p-1)}$. Its field population $A+Z$ is divisible
by $\operatorname{lcm}(8,2d)$. If denominator $p$ is absent, then
$A\ge d(p-1)$, $A+Z\ge d(p+1)$, and at least $dp$ nonunits have
pure $p$-power denominator. For $A<p(p-1)$ in this case $Z\ge2d$
and each moved neutral has $p^2\mid b_Y$. Always $A\ge2(p-1)$, and
$A\ge4(p-1)$ if $p\equiv1\pmod4$. At equality $A=2(p-1)$ the
active labels have a single denominator $B$ with $M_p=B^2/p$.
\end{lemma}
\begin{proof}
By localization, $p$-inertia fixes every label outside the
$p$-primary row fields, and these fields contain all active labels
since each row field contains its twist. On a $p$-primary nonunit
orbit the stabilizer in the cyclic local Galois group is odd, since
the field is nonreal. The generator of the local Sylow $2$-subgroup therefore acts in an
odd number of $d$-cycles. The active orbits
have size $(p-1)p^v$ since their stabilizers fix the square action on
a nontrivial twist and hence reduce to one modulo $p$. The determinant of this permutation on the
full label space is a sign. It is also a $p$-power root of unity,
since the local special linear group is generated by unipotents of
$p$-power order. Thus the number of moved orbits is even and
$2d\mid A+Z$.
The component column gives $8\mid A+Z$.

A phase of order $p^2$ costs $p(p-1)$ labels. At level $p$ the two
primitive square classes have equal mass per phase, so counts in pairs
agree modulo eight. Hence $2(p-1)\mid A$ and the number of neutral
moved orbits is even. An unrecorded prime of activity below $p(p-1)$
has an inactive pure label and no mixed denominator of $p$-part $p$
by the inactive-pure argument in Lemma~\ref{odd:inactive-pure}.
Its pure orbit costs $d$, and $Z\ge2d$. The neutral torus deficit is
at least $2(d-1)$, giving $A\ge d(p-1)$ by
Lemma~\ref{result:neutral-inertia}. The population modulus gives
the assertions also when $A\ge p(p-1)$. At $A=2(p-1)$ every primitive
phase is one pair of common degree by mass equality, and
$M_p=B^2/p$ by \eqref{eq:sparse-height}.
\end{proof}

\paragraph{The twist graph.}
A twist with three distinct prime divisors would reserve at least
$3\cdot5\cdot7=105$ duality orbits on denominators divisible by their
product, by Lemma~\ref{result:spherical}. For $r\le209$ this uses every
orbit and makes the mass at most $209/105^2<1$. Thus maximal supports
are vertices and edges, with
$\sum_{pt\ {\rm edge}}pt+\sum_{p\ {\rm isolated}}p\le(r+1)/2$.
Each $p$-active label at an isolated $p$ has pure $p$-power denominator.
At a prime denominator $p$ one has $p\Vert M$, $m_p=pj$ with
$a_p=(p-1)j$ active, and $p$ is isolated. Lemma~\ref{odd:field-component}
gives $p\ge11$. A prime denominator needs at least four support primes.
Its three proper degrees have gcd one, each involving two primes outside
$p$. The multiplicity $s=m_9$ is even and at most $76$ by
Lemma~\ref{result:saturation}. The remaining denominators have floor $15$.

At each recorded prime $p$ the proper $p$-free neutral labels occupy at
least three pairs at $3p$. With four primes these pairs may be chosen to
omit respectively each other prime $t$, giving divisors $pt^{v_t(M)}$.
If two prime classes occur their witnesses can overlap only in the pair
omitting both. On three-prime support every prime has height at least two
by the principal-block gcd argument. The nonzero-mean columns omitting
$p\ne3$ in nine-rows have three-primary twist and denominator divisible
by $9p^2$ for $s\ge9$, by Lemma~\ref{odd:nine-support} (by $3p^2$ otherwise).
Their numbers are bounded below by the paired version of
Lemma~\ref{result:joint-projection}, with lifts divisible by three
because every such degree contains three. Different omissions require
disjoint pairs.

\paragraph{Local moments at an edge.}
Fix a nontrivial complementary $t$-phase. Count dual pairs, let $p$
be the other endpoint and write $w_i=b_i^{-2}$, $V=\sum w_i$ for
the neutral pairs with $p\mid b_i$. Let $J$ be the active mass,
$C$ its mass at one primitive $p$-phase and $k$ the active count.
The integer Weyl coefficients $a_i\in p\ZZ$ have moments
\begin{equation}\label{odd:edge-mom}
 \sum a_iw_i=V-C,\qquad \sum a_i^2w_i=V+J .
\end{equation}
In a primitive fiber $J=(p-1)C$. Writing $a_i=p l_i$ gives
$\sum_i(pl_i-1)(l_i+1)w_i=0$. Since the odd weights equal one modulo
eight, the number $E$ of even $l_i$ is divisible by four. For $E=0$
every summand is nonnegative, hence all $a_i=-p$ and $C=(p+1)V$.
Denote odd-coefficient weight by $W_o$. The squares $(a_i\pm p)^2$
give $2p|V-C|\le V+J+p^2W_o$.

At a joint label the two local signs act independently. A product
of the complementary sign $C_{p'}$ with $p$-inertia cannot fix the
label because it inverts the nontrivial $t$-part of the twist in its
field (Lemma~\ref{odd:diagonal-twist}). In particular each fixed joint
phase and degree has an even number of dual pairs. By the Weyl parity
identity an even $a_i$ has odd stabilizer in $p$-inertia, so equal
coefficients at each degree occur in groups of $d=2^{v_2(p-1)}$ pairs.
$C_{p'}$ gives the disjoint translate. At level $p$, the $C_{p'}$-even
part of the fiber has active dimension $k$. The Borel-fixed projection
of $W_p^{-1}se_0$ supplies a Steinberg summand (its Weyl transform meets
every active phase). Conjugates of upper unipotents of level $p$ generate
the congruence kernel, as their leading terms span every graded
$\mathfrak{sl}_2(\mathbb F_p)$. The neutral part has at least $E(1-1/d)$
nontrivial torus dimensions, whence
$k\ge(p-1)(1+E(1-1/d)/2)$ by Lemma~\ref{result:neutral-inertia}.
For $p\equiv1\pmod4$ the eliminated moment modulo eight, grouped by
coefficients, even gives $8\mid E$.

For arithmetic with prescribed counts put $B=pt$ and scale weights by
$B^2$. Every larger mixed denominator is at least $3B$. A phase with
$c$ pairs, $u$ at $B$, has mass in $[u,u+(c-u)/9]$, strictly above $u$
if $c>u$. Apply this separately in both directions. Coefficients at $B$
are enumerated from their squared norm. The tail satisfies both moments
modulo eight. If $h$ neutral pairs have $v$ positions at $B$ of sum
$A$ and squared sum $H$, Cauchy on the tail gives
$(v-A-C)^2\le (h-v)(J-v+2A+2C-H)/9$. More generally one can sum
$(a_i-L)(a_i-L-p)\ge0$, $L\in p\ZZ$, on any tail, or use spacing
$2p$ on one parity and subtract $p^2$ times the opposite-parity weight.
These give the linear and quadratic interval tests for every fixed
small-denominator prefix. The bound on that weight is its count
divided by the squared next cofactor. A higher $p$-order costs
$p(p-1)/2$ pairs in addition to any primitive phases. Their positive
counts agree modulo eight by mass equality. At the largest denominator
$p$-height in a fixed fiber the normalized squared degrees sum to zero
modulo $p^2$ by \eqref{eq:phase-column}. A strictly higher level thus
requires at least three terms for $p\equiv3\pmod4$, and two otherwise.

\paragraph{Possible edges.}
Each edge $pt$ costs at least $4(p-1)(t-1)$ joint labels, by the
primitive-equal-mass or higher-orbit bound and the two independent signs.
For $p,t\ge5$ with $(p,t)\ne(5,7)$, this gives at least $160$ at
floor $55$ and at most $49$ others at nine, insufficient for mass one.
For $5,7$, a higher phase or an increase of two to four (or by eight in
one square orbit) in the common primitive phase count already gives
$192$ joint labels. Otherwise each joint phase has two pairs. At every fixed seven-phase
the second five-moment requires at least one solely-active pair. At every
fixed five-phase fewer than four solely-active pairs would force their
coefficients all equal to $-7$, giving $2\equiv8h\pmod8$ for their
count $h$. Thus at least $96+2\cdot6+8\cdot4=140$ labels have floor
$35$ and the upper mass $140/35^2+69/81$ is less than one.
For $3,t$ there are at least $16(t-1)$ $t$-active labels
of denominator divisible by $3t$. Each occupied nontrivial $t$-phase
has eight pairs by Lemma~\ref{phase:moments}. Thus $t\ge11$
gives mass at most $160/33^2+49/81<1$. Only $3,5$ and $3,7$ remain
to test, alone or together.

In a trial support with $2\sum p-1\le209$, enumerate the prime classes,
graph and $s$. Reserve the preceding disjoint groups. For a prescribed
number of slots at each floor $b$, their total mass is bounded by the
sum of the corresponding $b^{-2}$. At sparse isolated $p$ absent as
denominator the $p$-active degrees are $p$-divisible, so deletion
requires four neutral $p$-free labels at $p^3$. With four primes and
one prime denominator only one of its three proper pairs omits $p$.
On three-prime support, if $s\ge9$ then $27\mid M$ by
Lemma~\ref{odd:nine-support} and the principal gcd, so eight three-free
nonunits have floor $27$. The two sets of nine-row omissions and these
three-free labels are disjoint. All have three-primary twist, hence can
overlap mixed $3,t$ activity only when active solely at three. For each
possible overlap credit the larger of the two floors. Endpoints allow
only their own primes in active denominators by Lemma~\ref{result:spherical}.
For each allowed edge apply \eqref{odd:edge-mom} to the counts at
$B=3t$ and above. Square-unit symmetry identifies the two orbits of
third- or ninth-root phases and the square classes at $t$ separately.
At two leaves the $t$-active groups are disjoint. When a fixed fiber
has $n$ mixed pairs, at least three exceed $B$ if $3\nmid n$, and
one if $9\nmid n$, by the height test above. These rank and ordinary
mass checks leave only the empty graph and the edge $3,5$.

\subsection{The signed Gauss sum}
Average $\sum\theta_X/b_X^2=\xi/M$ ($\xi=\pm1$) over the cyclotomic
group. A label of twist order $e$ contributes to the deficit
$1-\xi/M$ the nonnegative rational weight
$(1-\mu(e)/\varphi(e))/b_X^2$ after averaging its whole orbit. Call
this weight its charge. At an occupied pure power $b=p^a$, the charge
is $pc/b^2$ with $c=h_1/(p-1)+h_2/p$, where $h_1,h_2$ count labels of
primitive and higher $p$-order. Thus $c$ is integral and costs between
$(p-1)c$ and $pc$ labels. A prime class contributes $a_p/[p(p-1)]$.
The nine class contributes $c/27$, $0\le c\le s/2$. When fewer than
$p^2(p-1)$ positions remain for the current and all higher levels, the height and balance bounds of Lemma~\ref{result:activity-height-gap} give even $c$. Every subsequent coefficient is below $p^2$, so the top charge has a unique largest pole at $p$. Its reduced denominator, and therefore all lower-level charge denominators, divide $M$.

Subtract the charges of complete specified groups from one, leaving $U$.
Let $L\mid M$ contain all known divisors. If the next denominator is $b$
and at most $k$ slots remain, then
\begin{equation}\label{eq:odd-deficit-tail}
 U\ne0,\qquad -1/\lcm(L,b)\le U\le1/\lcm(L,b)+3k/(2b^2).
\end{equation}
Indeed every remaining charge rate is at most $3/2$. If $U=0$, the
positive remaining charges would sum to $1/M$. Group them by their
component $J$. Their denominators are supported on $J$, up to a factor
two for the joint orbit of order fifteen. Comparing at the odd primes
of $J$ makes each positive group a multiple of $1/(2M_J)$, where $M_J$
is the $J$-part of $M$. At least two components occur, so $M/M_J\ge3$.
For an empty tail require $U=\pm1/M$. This argument gives finite
branching even before $M$ is known, since $\lcm(L,b)\ge b$.

\paragraph{The empty graph.}

Reserve six neutral nonunits overall (Corollary~\ref{result:mixed-neutral-rank})
or the stronger proper and omission counts above. Three-row omission
witnesses are neutral since three-activity has pure denominator.
The isolated field populations round up independently to multiples of
$\operatorname{lcm}(8,2d)$. With $s\ge9$ on three-prime support, three-height at least three
reserves eight three-free labels beyond the nine class, which may also
be active. For a mixed denominator $b=uk$ with $u$ its full $p$-part,
a dual pair spends at least $2\|k^{-1}\|_u^2/u^2$ in the local Weyl
norm (here $\|\cdot\|_u$ is the least absolute residue) (Lemma~\ref{result:local-entry-mass}). At a prime class use
its proper norm after reserving the other proper degrees. This bounds
the smallest possible free denominator, with the full tail retained
above the numerical cutoff.

Visit pure powers in increasing order. At $p$ let $t_p$ be the shortfall
in its activity quota and $B_p$ its next possible power. Include $B_p$
in $L$ when $t_p>0$, or $B_p^2/p$ when fewer than $2p(p-1)$ active
slots remain in that field. Indeed in the latter case the largest
active denominator $B$ satisfies $M_p=B^2/p$. The positive charge
still available is at most
$\sum_p t_p\alpha_p+(k-\sum_p t_p)\max_p\alpha_p$, where
$\alpha_p=p/((p-1)B_p^2)$. Allow error $1/L$ for the vacuum Gauss term.
Use the same reservation in ordinary mass. Enumerating $c$ at each
level, subtract its minimum label cost and credit its maximum activity,
rounding the field populations after each choice. \eqref{eq:odd-deficit-tail}
bounds every next level, and a terminal requires $L\mid M$.

\paragraph{The edge $3,5$.}
Let $z,x,y$ count respectively the joint, solely-three and solely-five
labels of mixed denominator. At primitive order fifteen, the eight
positive joint-phase counts in pairs are even and congruent modulo
eight, with a common count on each two-phase square orbit. An increase
of eight costs $32$ labels. A ninth-root square class costs at least
$48$ joint labels and has size divisible by $24$. Orders $25,27$ require
ten or nine phases containing at least four or eight joint pairs,
respectively. A trial with a higher order need only reserve one square orbit of
its phases, costing a multiple of $20$ or $18$ joint labels. Among the
surviving reservations the respective joint counts are $32,48,64,96,128$.
None is such a multiple. Thus the remaining phase counts are exact. A solely-active mixed label shares its phase with
joint activity, since $J=C=0$ in \eqref{odd:edge-mom} would give
$\sum(a_i-1)^2w_i=0$.

For $z=32,48$ there are four labels at each occupied full phase, by
the preceding counts. The sign orbits and equality of primitive
phase masses (square transitivity at ninth roots) give one joint
denominator $B$. In a three-phase there are eight joint pairs, two
at each primitive fifth root. The neutral torus bound gives $E<4$.
Thus every coefficient is $-5$ and the neutral pair mass is
$1/(3B^2)$. In units of $B^{-2}$ every pair has weight one modulo
eight, so its count is $3\pmod8$. Summing over two primitive or
three conjugate ninth-root phases shows $x\equiv12\pmod{16}$ or
$x\equiv18\pmod{48}$, respectively. For $J=z/B^2$ the solely-three
mass is $J/24$, primitive imbalance zero, and their combined charge
is $30/B^2$ or $50/B^2$. At $B=15$, three neutral pairs in a fiber
must all have denominator $45$ by their mass.

For a primitive fiber with zero-coefficient weight $Z$, the moments
give
\[
 \sum a_i(a_i-1)w_i=J+C,\qquad
 (V-C)^2\le(V-Z)(V+J),\qquad
 V\ge C/(p+1).
\]
A nonzero coefficient spends at least $p(p-1)w_i$ in the first sum.
At equality it is $p$ and all other coefficients vanish, giving
$V=C+pw_i$. Also $p|V-C|\le V+J$ bounds $V$ above. These facts
permit exact treatment of fibers where a position of denominator
fifteen would saturate the squared norm. Otherwise use the moment
intervals above, with next denominator at least $45$. In particular
in a primitive $32$-label fifth-root fiber a nonzero coefficient at
fifteen forces $V=5/225$. If there are fewer than thirteen solely-five
pairs all five must then have denominator fifteen.

\paragraph{Neutral positions.}
In the three-prime case keep disjoint three-free pairs, five-free
nine-row witnesses and $q$-free nine-row witnesses. At sparse isolated
$q$ their floors are $27,225,9q^3$. The first two can meet mixed
activity only in the solely-three sector, at floors $135,225$.
All $q$-free labels are neutral. A pure $q$-row makes their twists
$q$-primary, whereas $q$-active degrees are divisible by $q$.
To bound their number enumerate counts $v_t$ of pair residues
$t=1,\ldots,(q-1)/2$. Besides $1+2\sum v_tt^2\equiv0\pmod q$,
the quantity $2\sum v_t\|at\|_q^2$ must reach the projection bound
\eqref{eq:full-census} for each $a$, using the two consecutive
multiples of three in $-a+q\ZZ$ and even selected-row multiplicities.
Additional $q$-free labels have floor $q^3$.

For four primes with prime denominator $p$ and other isolated prime
$q$, proper pairs omit $3,5,q$ at floors $9p,5p,pq^3$ in the sparse
case. Only the last may be $q$-free. The seven-free count is at least
six and at least eight when the mass of the nine and marked classes
exceeds $3/7$, by the same residue bound. For $q\ne7$ reserve an
additional pair. With two prime classes $p,q$, the five distinct pair
floors $9p,5p,9q,5q,pq$ suffice. All independent reservations are
combined by maximizing mass over their possible overlaps.

\paragraph{The nine-class projection.}
Let $F$ be the nine class of size $s$, $I$ its support together with $F$,
$H=\sum_I b_Y^{-2}$, $E=H-s/81$, $C_l=\sum_F\theta_Y^l/81$.
Let $v=(se_0)|_I$, $P$ the orthogonal projection onto the rows $se_X$,
$X\in F$. Then $Pv=0$, and the coordinates of $PT^{\pm1}v$ on these
rows are $\epsilon\theta_X^{\mp1}/9$ by the modular relation. Pairing
$(1-P)Tv$ with $v$ and $(1-P)T^{-1}v$ gives
$(|C_1|-E)_+^2\le HE$, $|C_2-\bar C_2|\le2E$. For nine-charge $c/27$
and absolute primitive imbalance $i/81$,
$81|C_1|=\sqrt{(s-3c)^2+3i^2}$ and
$81|C_2-\bar C_2|=2\sqrt3 i$ after averaging higher-order phases
over their wild Galois subgroup when necessary. Indeed at these
complete levels their primitive-root sums vanish also before averaging.

For $A=(T+T^{-1})/2$ the squared residual norm of $Av$ is the
weighted sum of squared real parts on $I\setminus F$, since the
$F$-contribution cancels the projection. The Galois means of a real
part and of its square are $t$ and $(1+t)/2$, $-1/2\le t\le1$, by
odd order. If $e_0,u_0,h_0$ are the specified tail mass and its first
two real moments, with unspecified mass at most $E'$, Cauchy and
averaging give
\[
 2((s-3c)/81+u_0+u)^2\le(s/81+e_0+E')(2h_0+E'+u),
 \qquad -E'/2\le u\le E'.
\]
At a nine-multiple $b$, $a$ solely-three labels with $i'$ primitive
packets add $(a,-i',(a-i')/2)/b^2$ to the three specified moments.
Every supporting label has three-primary twist and is at a multiple
of nine by Lemma~\ref{odd:nine-support} in the cases tested. Use the
least common multiple with nine of each omission divisor.

\paragraph{The mixed recursion.}
Enumerate the occupied denominators in increasing order, assigning
complete square-unit packets and their duals. The charge per label is
$7/(8b^2)$ at order fifteen, $1/b^2$ at $45$, $p/((p-1)b^2)$ at order
$p$ and $1/b^2$ at higher prime powers. For $p=3,5$, if $i_p$ primitive
packets occur, update the absolute difference $d_p$ of the two
primitive phase masses per phase at trivial complement to
$|d_p+j/b^2|$, $|j|\le i_p$, $j\equiv i_p\pmod2$.
Start with $d_3=i/81$, $0\le i\le c$, $i\equiv c\pmod2$, and $d_5=0$.
Both vanish at completion. Subtract the minimum label cost and test
all residual ordinary, charge, phase and Gram intervals. In
\eqref{eq:odd-deficit-tail} one may use $M\ge\kappa b$, with $\kappa$
the least of $L/(L_3L_5)$ and $pL/L_p$ at sparse isolated primes
($b^2/p\mid M$ for a pure active $b\ge p^2$). With next level $b>15$
a free label contributes at most $1/225$ to ordinary mass plus
$2(b^2-225)/675$ times its charge. For remaining solely-three,
joint and pure-three masses $A,J,P'$ with $P'\le P$, booked
$A_0,J_0$ and $A+J+P'\le T$, the fiber bound
$16(A+A_0)\le9(J+J_0)$ and $2d_3\le A+P'$ imply
$50d_3\le9T+9J_0-16A_0+16P$. At order $45$, $2d_3\le P$.
At minimum isolated activity reserve the neutral pair at the unique
active denominator by Lemma~\ref{result:packet-idempotent}.
Otherwise use Lemma~\ref{odd:inactive-pure} for its floor.

A terminal determines $M$ by $U=\pm1/M$. Check divisibility by every
specified denominator and the prime heights. The residual ordinary
squared-degree sum is twice a sum of squares of admissible divisors
$d$ of $M$ at floor $M/d\ge15$, with the number of pairs congruent to
that half-sum modulo eight. Enumerating multiplicities of the six
largest available squares and bounding the rest by their floors gives
a necessary interval test. Bounds on activity and residual mass always
use the ceiling $209$. Any count known only to be below the ceiling
retains the whole range.

\paragraph{Small mixed fibers.}
We record the integer consequences of the moment bound used at the last
levels. Scale by $15^2$, set $u_i=a_i-1$ and let $k_i=b_i/15$.
A fiber of $j$ joint pairs satisfies
$\sum u_i/k_i^2=-C$, $\sum u_i^2/k_i^2=J+2C$, $E\equiv-j\pmod8$,
the last by summing squares modulo eight in the even coefficient groups.
At any of the first cofactors $1,3,5,9$, enumerate
$u_i^2\le k_i^2(J+2C)$. On the unassigned tail use the moment and parity
intervals above. Modulo $32$ each equal-cofactor group uses one of
$1,9,17,25$ as common weight. If its mass is fixed to $v>0$, the next
cofactor is at most $\sqrt{h/v}$ for $h$ remaining pairs.

At $z=64$ primitive, divide each phase into sign orbits of four labels.
The positive numbers of such orbits per phase are equal modulo eight,
constant on two-phase orbits, and total sixteen, hence all two. If both
groups in a phase have the smallest denominator, the same holds in
every phase. If exactly one does in each phase, their second denominators
agree. A phase can omit the smallest $b$ only if the next $b'$ has
$b'^2<2b^2$. For $z=96,128$ with all joint denominators fifteen, and for each feasible decomposition into $15$- and $45$-order phases,
$z=16a+24(b+c)$ with $a,b,c$ even, the local $(J,C)$ are $(4t,t)$
at the three-phase of type $t=a,b,c$ and $(2a+3b+3c,a)$ at each fifth
root. These follow by counting the full phase pairs, with $a$ at each
primitive phase and $b,c$ on the ninth-root square classes.

In each coefficient group the nonnegative moment
$\sum a_i(a_i-1)/k_i^2=J+C$ charges at least $p(p-1)/k_i^2$ per
odd coefficient, and at least $2p(2p-1)/k_i^2$ per nonzero even one.
If the odd coefficients exhaust this budget at their minimum, they
all equal $p$, the even coefficients vanish, and the first moment fixes
the mass. For groups assigned cofactors below $75$ use their exact
weights and bound each remaining weight by $1/75^2$.

Below $84$ active labels at seven, the possible counts are
$12u+42v$, $v=0,1$, by the square classes of primitive and higher twists.
Reserve the groups at $49$, then $343$, bounding the rest from $2401$.
A lone final group has charge $14/B^2$ or $42/B^2$, with $B^2/7\mid M$.
In the sparse case the total count is $12,24,36$. A pure isolated row
on three-prime support requires separate three-free and five-free neutral
pairs at floors $3^{H_3}q^2,5^{H_5}q^2$, where $q$ is the isolated
prime and $H_i$ the known height, by principal support and gcd one.
They can coincide with the corresponding omissions already counted.
Initially a free label of floors $15$ neutral or $25,27$ active has
charge at most $1/486$ and ordinary mass plus $64/45$ times its charge
at most $1/225$.

To check remaining slots modulo eight, each primitive or higher packet
at three adds respectively (charge numerator, count) $(3,2),(6,6)$.
At five the pairs are $(5,4),(20,20)$. Odd reciprocal squares reduce
to one modulo eight. Keep the charge residue and the two primitive
imbalance parities, admitting every odd residue of $\pm1/M$. Shortest
path length in label count gives a lower bound. At the largest sparse
seven-power $B$, clearing denominators gives $M_7=B^2/7$ and
$cM/M_7\equiv-\xi\pmod{49}$, $c=2,4,6$. A prime class at eleven gives
$2M/11\equiv-\xi\pmod{11}$ and hence $\xi=-1$ on the support
$3,5,7,11$ when $M_7=B^2/7$. The seven-height is odd, $3,5$ are
quadratic residues and $2,7$ nonresidues modulo eleven. Testing
multipliers of $L$ up to $10^5$ gives a lower bound for $M$. The
entire larger tail is retained.

At ceiling $209$, the prime support and prime-class enumeration has
$16\,512$ profiles. The graph and nine-class tests leave $1\,151$ edgeless
reservations and $310$ with the single edge $3$--$5$. The two recursions
visit $6\,721\,348$ and $2\,753\,944$ nodes respectively and leave no admissible completion. Since $r\equiv1\pmod8$, the next odd
square rank is $217$, proving the odd-dimensional assertion of
Theorem~\ref{alg:rank}.

\section{The denominator census}\label{sec:verify}

Suppose $r\le44$. By Section~\ref{sec:odd-new}, $D$ is even.
Let $B$ be the least nonunit denominator, of multiplicity $m$, and put
$n=M/B$, $f=qB^2$. The ordered denominators determine $M$, their lcm,
since the nonunit degrees have gcd one. We have $f<r$. Equality
$f=r-1$ would give $\sum_{i\ne0}(n^2-d_i^2)=1$, whereas each
positive summand is at least $2n-1$. The small-codegree bounds thus leave
$16,25,36$ for square $D$ and $12,18,20,24,27,28,32,40$ otherwise.
Apply Lemmas~\ref{result:small-effective} and~\ref{block:moments}, and
$m\le f-1-(p-1)/2$ for odd $p\Vert f$
\cite[Proposition~19.9]{BDPcore}, by principal height and the primary mass identities.

\paragraph{The search interval.}
Use \eqref{eq:interval}, with $h$ nonunit positions still empty.
For the two-degree terminal equation it is enough to try
$k^2\le(1+2L^2)/E$ and the divisors of the prime-to-$k$ part of $L$.
With three degrees $d,e,g$ left and $k$ even,
$Ek^2=1+d^2+e^2+g^2$ modulo four and eight makes all three odd,
$v_2(k)=1$, and $E$ odd. No new prime can enter $M$ at the last
two positions, by Lemma~\ref{block:witness}. 

\paragraph{Reserving positions.}
If $v$ of the $h$ missing nonunits have respective denominator floors
$a_i$ and all denominators are at least $b$, the residual mass is at most
$\sum a_i^{-2}+(h-v)b^{-2}+M^{-2}$. Combine overlapping reservations
by taking coordinatewise maxima after sorting. For
$W=\sum\lceil L^2/a_i^2\rceil$ and $E>W+1$ this gives
$b^2\le(h-v)L^2/(E-W-1)$, since $L\mid M$. The activity bound on
every set of odd primes $p\mid M$, $p\nmid q$, lets Hall's theorem reserve
disjoint groups of $B(q)(p-1)$ labels for these primes. A closed class of
size $s$ supplies at most $h_p\lfloor s/h_p\rfloor$ at $p$,
$h_p=(p-1)/2$. Credit an open class in full. The support contains the
primes of $2qB$ and satisfies $B(q)(1+\sum_{p\mid M,\,p\nmid2q}(p-1))\le r$.

For $q=3$ the neutral fiber at the unramified odd primes costs four
labels in place of $B(3)=3$. Its principal ternary block contains the
vacuum. A phase of order at least nine supplies three more labels.
Otherwise the order-three mass is twice the inactive mass. With at most
three labels, all nonunit degrees exceeding two, this requires
$b^2=2(1+a^2)$ for nonunit degrees $a,b$. Modulo three they are
three-free and modulo four $a$ is odd, so $(a,b)=1$. Their $S$-entry
vanishes, contrary to principality at three.

For $q=3,5$, $p\ge5$ distinct from $q$, put $F=(q+3)/2$.
Every nonempty fiber away from $q$ costs $F$ labels, even with dyadic
phase fixed, as the complementary projectors commute with the $q$-factor.
With fewer than $2F(p-1)$ active at $p$, their phases have order $p$:
order at least $p^2$ costs $Fph_p$ labels. The fiber with trivial
complementary phases contains an active $p$-phase, since otherwise the
local unipotents would preserve the vacuum line. Equal primitive masses
give all $p-1$ roots there. If an active $X$ has $t\mid f_X$ outside
$\{p,q\}$, its fiber away from $t$ must contain two $t$-phases by the
same invariant-line argument. Thus some nontrivial $t$-phase also contains
all $p-1$ roots, costing $2F(p-1)$ labels. Hence active denominators are
supported on $p,q$. At $q$-free denominators their defect-one blocks must
have equal degrees, since a mixed block includes $z,2z$. Their activity
multiplicity is consequently divisible by $qh_p$.

Tile $p$-free denominators for $p\mid q$. At two the counts are even.
At three and five, after obligations from $z/2$, a mixed block uses
respectively one or two copies of $z$ and two copies of $2z$.
The least mixed count permitting equal blocks at $z$ suffices: $p$
mixed blocks may be exchanged for equal blocks at both levels.

At $q=3$ on support $\{2,3,p\}$, any position at denominator $2$ or
$p$ needs two at twice its denominator. An equal ternary block there
would require a fourth prime by Proposition~\ref{result:effective}.
At degree $d$, constituents of $XX^*-1$ prime to each $p\mid d$
have total degree congruent to $-1$ modulo the least $p$-part of a
$p$-divisible degree. On three-prime support a nonunit can omit only
one prime. Thus for the three least nonunit degrees $d\le e\le a$,
with repetition, $e+a\le d^2-1$, since at least two distinct
constituents other than $X$ are required.

For $p\ne t$ in the support, a nonunit in $P_p$ omits $t$, giving a
denominator multiple of $pt^{v_t(M)}$. Orbit deletion forces two
$t$-free nonunits, and reduction of $D$ modulo four or five forces
three for $t=2,5$. With three support primes their omissions are
disjoint, every $t^2\mid M$, and the two- and three-free counts are
$3\bmod4$ and $2\bmod3$, respectively. At an undecided height take
the weaker bound for a fixed or increased valuation. With four
primes a nonunit omits at most two. Distinct omitted pairs require
distinct positions. The proper support at a prime denominator must
omit each other prime somewhere by its gcd-one condition.

\paragraph{Square dimension.}
There are four primes, and the only possible edges are $2$--$3$ and
$2$--$5$ (Lemma~\ref{result.Three-prime-rank} and the subsequent support
bounds). Use the multiplicity restrictions of Proposition~\ref{result:thirteen-34}. At an isolated odd $p$ reserve its $p$
pure nonunit positions. At $p\ge7$ there are at most $37$ active labels
after reserving the vacuum and the two other odd activities. At seven
this number is a multiple of three, hence at most $36$. At isolated
five with no denominator five there are at most sixteen. For at least
nine fours this is the high-four bound, and with at most eight fours
eighteen active labels would give total mass at most
\[
 18/625+6/49+2/36+8/16+9/36+1/210^2<1 .
\]
The count is even. Thus at every isolated $p\ge5$ the sparse-height
inequality applies. In the absence of denominator $p$, no mixed
denominator has $p$-part $p$.
At the greatest active level $p^a$, $v_p(M)=2a-1$ and at least
$p-1$ labels occur there (a higher-order packet suffices, and for
primitive packets reduce equality of the phase masses modulo $p^2$
after scaling by $p^{2a}$). These give denominator floors both for
the active labels and the omitted-prime witnesses. At isolated three reserve only two equal pure denominators, by the same
primitive-mass reduction or a higher-order orbit.
If seven occurs, five does not, and $m_4+m_6+27>r$, five cannot be
mixed: it would need sixteen active labels in addition to the vacuum,
the fours and sixes, and either fourteen sevens or seven sevens and
three proper, twist-trivial witnesses.

For an isolated $p$, let $E_p$ be the mass of all pure $p$-power
nonunits and $C$ the sum, over observed mixed $b=uk$ with $u$ its
$p$-part, of $\min_{ak\equiv1\, (u)}(a/u)^2$ ($u>1$). Since the pure
entries use $E_p$ of the squared norm $2g_p-g_p^2$ and
$0<g_p\le pE_p/(p-1)$,
\[
 C<(p+1)E_p/(p-1).
\]
With fewer than $p^2$ active labels, equality of the primitive masses
at successive levels gives equal base-$p^2$ digits at each pure
denominator. Higher-order orbits contribute in multiples of $p$
and have size at least $ph_p$. Thus each class of $n_b$ labels contributes
$pc_b/b^2$ to $g_p$ with $(p-1)c_b\le n_b$. For $H$ remaining
pure positions of floor $b_*$ it follows that
\[
 C\le\sum_{b\ {\rm closed\ pure}}
 \frac{2p\lfloor n_b/(p-1)\rfloor-n_b}{b^2}
 +\frac{(p+1)\lfloor H/(p-1)\rfloor}{b_*^2}.
\]
Indeed $2pc_b$ norm units use at least $(p-1)c_b$ units on the pure
positions themselves. If a mixed dyadic--$p$ phase occurs,
each of its required fibers uses four positions of one positive dyadic
height. In the ternary case the two primitive masses agree. A partial
difference at denominators of lcm $W$ is a multiple of $1/W^2$. If the
remaining eligible mass is smaller, the difference is zero.
Between two such cuts a singleton at maximal dyadic or ternary
denominator valuation must have coefficient zero, by reduction at that
prime. A singleton also cannot carry an order-nine phase. Removing these
positions before reserving the quartets preserves every solution.

At a prime denominator $p$ of multiplicity $pj$, use the signed
coefficient equations of Proposition~\ref{result.One-packet} on its proper
witnesses. For $p=7$ the budget
$j(7-j)\le12$ forces each coefficient to be the least absolute
inverse of $b/7$ modulo seven. For any partial list, with remaining
budget $K$, its missing sum $\sum a/(b/p)$ is bounded in either
sign by $K/\min |a|(b/p)$, minimizing over admissible remaining
terms of that sign. With one packet at $p=5,7$ write $W$ for the
lcm of a nonempty proper initial witness list, $s=\sum aW/b$.
Then $s\equiv (W/p)\sum a^2\not\equiv0\pmod p$ since the total
budget is $p-1$. For the next denominator $z$, if the missing
absolute coefficients sum to $K$, the signed equation gives
$z\le(K+1/6)W/|s|$. At the last witness put $d=W/\gcd(W,z)$,
$t=z/\gcd(W,z)$. The gcd-one condition gives $M=Wt$ and
$st+a_zd=\pm1$. Testing $d\mid W$ is exhaustive.
For simultaneous single packets at five and seven at most two proper
labels are shared. In each family separate labels must omit each of
the two other primes. In the three-label seven family at most one
can be shared. The cross-column sum of products on the shared
labels is $-1$, so exactly one is shared, with opposite signs,
also when the seven coefficients are all signs. Its denominator
$35k$ has $k\equiv\pm17\pmod{35}$. The remaining seven witnesses
are obtained from their same signed equation, with the common sign
specified by the vacuum congruence. With a present five, a mixed
ternary phase and $m_4\ge9$, the large isolated prime has at most
$r-m_4-20$ pure positions, after the vacuum, fours, five-packet,
four proper witnesses, two further odd degrees and eight mixed positions.
When this is at most $2(p-1)$ and $p$ itself is absent, the minimal
packet bound gives $p$ positions at the first pure level $B=p^c$ and
$M_p=B^2/p$, with $M/M_p\equiv-\epsilon\pmod{M_p}$.

\paragraph{Gauss sums.}
For a complete class write $T_b=\sum_{b_i=b}\theta_i$.
If selected classes have denominator lcm $W$ and the complementary
mass is at most $R$ after discarding known zero sums, every conjugate
of $W^2\sum T_b/b^2$ has modulus at most $W^2(R+1/M)$.
If this bound is below one, integrality forces the sum to vanish. Reduction modulo
a prime detects nonvanishing when a unique maximal-height class has
one label, or is a pure $p$-class with multiplicity prime to $p$.
One may repeat after each zero sum, or apply the dyadic cancellation
bound to two classes. A selection of individual terms suffices as well. After
a zero seven-sum, add one twelfth-denominator twist to $9T_4+4T_6$
when $m_4$ is even. Reduction at two then leaves a single unit.
When $m_4$ is odd the four-class alone gives that unit. This proves the corresponding suffix tests by
using an upper bound for $R$. For example, with ten fours, five fives
and six sixes the five-sum vanishes. The short Gauss relation with
$72(1/120+1/840)<1$ gives $T_4=T_6=0$. The remaining mass is
$1/120$. A fourth prime $p\ge7$ is isolated and absent as a denominator
since $p-1\le41$ and a prime packet alone would exceed that mass.
Thus $M\ge8\cdot3\cdot5\cdot7^3$.

For an isolated odd prime $p$, put
$M_p=p^{v_p(M)}$. After subtracting the deficits
$(1-\langle\theta_i\rangle)/b_i^2$ of any complete nonunit classes
from one, the remainder is nonzero. Otherwise the unassigned
nonnegative deficits sum to $1/M$ by the Gauss identity.
Their $p$-active part $P$ has $2P$ of $p$-power denominator. Only
pure $p$-denominators contribute, and every class sum is invariant
under square multipliers, whose index is two at $p$. All other
nonzero terms have denominator prime to $p$, since all classes are
square-multiplier invariant and their Galois means have
power-of-two denominator. Hence the denominator of $2P$ divides
$M_p$, giving either $P=0$ or $P\ge1/(2M_p)>1/M$.
The former contradicts the denominator of $1/M$. The unassigned
deficits are bounded by twice their mass, with terminal error $1/M$.
At isolated $p$ with fewer than $p^2$ active labels one uses
the digits $pc_b$ above. At all other classes the root means are
$\mu(k)/\varphi(k)$.

\paragraph{Excluding a complete list.}
Apply five tests in order. First, if $d$ is the least nonunit degree and
$E_p=\sum_{i\ne0,\,p\mid f_i}d_i^2$, Lemma~\ref{block:moments} requires
$2pdE_p\ge(p-1)D$ for every $p\mid D$.

Second, at a prime denominator $p\nmid q$, put $n=M/p$. For $q=1$ this
gives a row of codegree $p^2$. For $q=2$, or for $q=3,5$ when neither
$2n$ nor $n/2$ occurs, the defect-one block of the ramified prime has
equal degrees and gives an effective row of codegree $p^2$.
Proposition~\ref{result:effective} then requires gcd one
among the nonunit degrees prime to $p$.
Third, tile the defect-one classes as above.

Fourth, let $A_p$ consist of the nonunit degrees $y$ with $p\mid f_y$ such
that, for every $p$-free nonunit degree $x$ of multiplicity $m$,
\[
 g=(x,y)>1,\qquad 2mg>p,\qquad m\,\lcm(x,y)^2<D .
\]
It contains every nonunit degree of $P_p$, so $\gcd A_p=1$.
Indeed a principal row meets the entire $x$-class, giving $g>1$.
Its entries divided by $x/g$ are nonzero algebraic integers.
Their sum is an integer congruent to $mg$ modulo $p$ and strictly
smaller than $mg$ in modulus. Their mean squared moduli are at
least one. The vacuum uses one further unit of row norm.

The remaining lists have $q=1$ and a denominator five. By
Lemma~\ref{result:neutral-inertia} there are four proper five-free
nonunit labels, i.e. those of degree $d\ne M/5$ prime to five.
On these an active five-row has signs $s_d\equiv d\pmod5$
by \eqref{eq:packet-trace} and principality. It vanishes on
five-divisible degrees. Its own degree class lies outside $P_3$
by principal height. Thus some subset of the four labels with
$d\in A_3$ must satisfy $1+\sum s_d d\equiv0\pmod{M_3}$
by \eqref{eq:principal-trace}.

\begin{proof}[Conclusion of the proof of Theorem~\ref{alg:rank}]
The denominator recurrence gives the following numbers of complete lists.
\[
\begin{array}{c|rrrrrrr|r}
 r&1\text{--}38&39&40&41&42&43&44&\text{total}\\ \hline
 q=1&0&1&2&17&111&914&6269&7314\\
 q>1&0&0&0&4&65&723&10093&10885
\end{array}
\]
Of these $18\,199$ lists, the five tests leave successively
$8\,207$, $6\,194$, $6\,166$, $244$ and $0$. The enumeration covers
$12\,273$ tasks. At a fixed-width arithmetic limit the unfinished branch
continues in arbitrary precision. There are $41\,915$ such continuations.
The supplement records each list and its first exclusion. Running
\texttt{python3 verify.py --full} regenerates the census. Every admissible
list follows the ordered recurrence and satisfies its pruning inequalities.
Hence $r\ge45$. Section~\ref{sec:odd-new} gives the odd-dimensional bound.
\end{proof}

\subsection*{AI use declaration}
GPT-6 Astra and GPT-6.2 Astra assisted with the development and drafting of proofs, the literature search, and the implementation and checking of the finite calculations.

\begin{thebibliography}{99}
\bibitem{ABPP}
M.~A.~Alekseyev, W.~Bruns, S.~Palcoux and F.~V.~Petrov,
\href{https://doi.org/10.1063/5.0334158}{\emph{Classification of integral
modular data up to rank 13}}, J. Math. Phys. \textbf{67} (2026), no.~8, 081701. Numbering follows the \href{https://arxiv.org/abs/2302.01613v10}{open author version}.

\bibitem{BDPcore}
S.~Burciu, J.~Dong and S.~Palcoux,
\href{https://hexagonmath.org/2611.00233v3}{\emph{Prime factors of fusion-category dimensions: weak group-theoreticity, rank and decomposition}},
Hexagon preprint 2611.00233v3 (2026).

\bibitem{BPBurnside}
S.~Burciu and S.~Palcoux,
\href{https://doi.org/10.1142/S0219199726500070}{\emph{Burnside-type
results for fusion categories}}, Commun. Contemp. Math. (2025), article 2650007.

\bibitem{CMSCyclotomic}
F.~Calegari, S.~Morrison and N.~Snyder,
\href{https://doi.org/10.1007/s00220-010-1136-2}{\emph{Cyclotomic integers, fusion categories, and subfactors}},
Comm. Math. Phys. \textbf{303} (2011), 845--896.

\bibitem{DLNCongruence}
C.~Dong, X.~Lin, and S.-H.~Ng,
\href{https://doi.org/10.2140/ant.2015.9.2121}{\emph{Congruence property in conformal field theory}},
Algebra \& Number Theory \textbf{9} (2015), no.~9, 2121--2166.

\bibitem{ENOannals}
P.~Etingof, D.~Nikshych, and V.~Ostrik,
\href{https://doi.org/10.4007/annals.2005.162.581}{\emph{On fusion categories}},
Ann. of Math. \textbf{162} (2005), no.~2, 581--642.

\bibitem{HR}
S.-M.~Hong and E.~C.~Rowell,
\href{https://doi.org/10.1016/j.jalgebra.2009.11.044}{\emph{On the classification of the Grothendieck rings of non-self-dual modular categories}},
J. Algebra \textbf{324} (2010), no.~5, 1000--1015.

\bibitem{LPR}
Z.~Liu, S.~Palcoux, and Y.~Ren,
\href{https://doi.org/10.1142/S0129167X23500301}{\emph{Interpolated family of non-group-like simple integral fusion rings of Lie type}},
Internat. J. Math. \textbf{34} (2023), no.~6, 2350030.

\bibitem{Nevins}
M.~Nevins, \href{https://doi.org/10.4153/CJM-2005-026-1}{\emph{Branching
rules for principal series representations of $\mathrm{SL}(2)$ over a
$p$-adic field}}, Canad. J. Math. \textbf{57} (2005), 648--672.

\bibitem{NSind}
S.-H. Ng and P. Schauenburg, \href{https://doi.org/10.1007/s00220-010-1096-6}{\emph{Congruence subgroups and generalized Frobenius--Schur indicators}}, Comm. Math. Phys. \textbf{300} (2010), 1--46.

\bibitem{PrasadSL}
A.~Prasad, \href{https://arxiv.org/abs/0712.4051}{\emph{Representations
of $\mathrm{GL}_2(\mathbb F_q)$ and $\mathrm{SL}_2(\mathbb F_q)$, and
some remarks about $\mathrm{GL}_n(\mathbb F_q)$}}, 2007.

\bibitem{Stacks}
The Stacks Project Authors, \emph{The Stacks Project},
\href{https://stacks.math.columbia.edu/tag/04GH}{Tag 04GH}.
\end{thebibliography}
\end{document}
